\documentclass[12pt,letterpaper]{article}
\usepackage[utf8]{inputenc}
\usepackage{amsmath, amsfonts, amssymb, amsthm}
\usepackage[top=1in, bottom=1in, left=1in, right=1in]{geometry}
\usepackage[backend=bibtex]{biblatex}
\usepackage{mathdots, mathrsfs}
\usepackage{titlesec, hyperref, textcomp, parskip}
\usepackage[]{footmisc}
\usepackage{graphicx, svg}
\usepackage{algpseudocode}
\usepackage{comment}

\titleformat{\section}{\normalsize\scshape\center}{\thesection}{1em}{}
\titleformat{\subsection}{\normalsize\scshape\center}{\thesubsection}{1em}{}

\newtheorem*{definition}{Definition}
\newtheorem*{theorem}{Theorem}
\newtheorem*{proposition}{Proposition}
\newtheorem*{lemma}{Lemma}
\newtheorem*{conjecture}{Conjecture}
\newtheorem*{observation}{Observation}

\newcommand{\ext}{\textnormal{ext}}
\newcommand{\B}{\textnormal{B}}
\newcommand{\mb}{\mathbf}
\newcommand{\mc}{\mathcal}
\newcommand{\sign}{\textnormal{sign}}
\DeclareMathOperator*{\argmin}{arg\,min}
\DeclareMathOperator*{\argmax}{arg\,max}

\begin{document}

\title{\normalsize{\uppercase{\textbf{Notes on learning, intelligence, and open-endedness}}} \\ (mainly open-endedness)}
\author{\small{\textsc{Lukas Nabergall}}}
\date{\small{\textsc{\today}}}
\maketitle

\allowdisplaybreaks

\textbf{Ackowledgements}: Much of the following is based on the ideas of Kenneth Stanley, Joel Lehman, and other researchers in the fields of artificial intelligence and open-endedness.

\textbf{Some background motivation}: To the best of our knowledge, the only known ``process" to ever produce intelligence is (biological) evolution. We put the descriptor {\em process} in quotes because its perhaps somewhat of a restrictive conceptualization to think of evolution in such monolothic terms, rather than simply as a type of persistent physical phenomena. Before human technology, evolution generated all known examples of intelligence, namely all living organisms. It also generated us, the only known type of system endowed with what is usually termed {\em general intelligence} (or, a bit less vaguely, human-level intelligence).\footnote{More could be discussed regarding this notion of general intelligence.} More broadly, evolution is arguably the origin of most of the novel phenomenon that we observe. The universe is a complex and interesting place mainly due to the complexity and diversity of life and its innumerable effects on Earth, and evolution continues to generate novel things at a rate that seems to far exceed the rest of the known universe. Note that its important to distinguish novelty from information. For example, while a star is a high information object relative to the thin biosphere of Earth, any given main sequence star is highly similar in most interesting or useful ways to many other stars, so there is little novelty to be gained in observing the $n$th main sequence star after observing $n-1$ previous main sequence stars.\footnote{Much more also could be (and should be?) discussed regarding the notion of novelty.} In contrast, the $n$th species is often highly novel relative to a sampling of $n-1$ previous species, exhibiting new behaviors, physical structure, etc.\footnote{This is also potentially true for individual organisms, but at least less obviously so.}

These features of evolution are captured by the concept of {\em open-endedness}. A system or process is {\em open-ended} if it produces abundant novelty indefinitely, seemingly or practically without end. Even after a very large amount of time has elapsed, perhaps millions or billions of years, the system continues to spit out novel artifacts, take on novel features, undergo novel events. Such a system exhibits increasing complexity and endless creativity. Like most things in these notes this should not be taken as an attempt at any sort of complete, explicit definition but just a description of what we are getting at when we use the term `open-endedness'. Any attempt to further understand this description should usually appeal to the systems that embody and inspire this open-endedness property.

%%% maybe should think a bit---is the universe an open-ended system? And should this not just be set aside, but commented on and considered?
%%% is it open-ended only because of the evolution of life, humans, etc. or are those just an extension of its prior open-endedness? Life could be the continuation of its novelty, which would otherwise have gradually dried up....
%%% this also potentially gets at the distinction between thinking about evolution at what is usually viewed as the evolution-level (DNA, genes, cells, etc.) and thinking about evolution at the "universe-level" (energy and matter differences spark a series of chemical reactions that eventually lead to the jumpstart of evolution, etc.)

There are arguably essentially only three known open-ended systems: 1) the observable universe, 2) the biosphere, generated by evolution, and 3) human civilization. We distinguish these systems by their underlying generative processes, despite the fact that there is significant overlap between them---all observable things including the biosphere are by definition part of the observable universe, while humans are part of both the biosphere and human civilization. We should note though that system 1 holds a unique position in comparison to systems 2 and 3.

The observable universe is everything that we know. Unlike the biosphere and human civilization, it is not commonly understood to be an open-ended system.\footnote{I think?} Nevertheless it seems both reasonable and important to make this classification. The universe generated both the biosphere and human civilization and the dynamics that drive those open-ended processes do appear to be the same dynamics that govern the universe as a whole. Furthermore, the universe contains a large, diverse collection of objects that grew and varied throughout its existence, including stars, planets, comets, asteroids, nebulae, galaxies, black holes, gas clouds, brown dwarfs, asteroid belts, galaxy clusters and superclusters, globular clusters, quasars, etc.\footnote{In no particular order...} Each of these types of objects come in a great variety of forms exhibiting an immense plethora of phenomena, most notably amongst the planets. Although the current leading theory forecasts that the universe will end in the decidedly non-open-ended heat death, this arguably no more mitigates the case for its open-endedness than it does for the open-endedness of the biosphere or human civilization, both of which will\footnote{As far as we know.} suffer the same fate. It is nevertheless worth acknowledging that there are seemingly substantial differences between system 1 and systems 2 and 3. Although we will not dwell on elucidating these differences, which are far from completely apparent, the most obvious distinction is that systems 2 and 3 are driven by (approximate) replication mechanisms, namely reproduction of living organisms and humans, respectively, as well as the reuse and adaptation of ideas, tools, and structures in human civilization. In contrast, replication is not a {\em central} mechanism in the dynamics of the universe, at least not on most scales.

%%% discuss that these are open-ended, and why they are the only examples
%%% every other system is a subprocess, especially in so far as they are also directly generated by humans (e.g. the internet); the key is separation of the generative mechanisms
%%% why open-endedness? two examples, etc.

While all three systems deserve consideration, we think there are a number of reasons to focus on the universe (qua universe) in our initial open-endedness studies and briefly describe them here:
\begin{itemize}
\item The universe is the parent system in which all open-ended systems we know of are contained. All other open-ended processes are generated by the same dynamics that generate the universe, even if higher-level abstractions are a more useful or apt explanatory language. This asymmetry suggests that studying the universe will provide more information on other open-ended systems than vice versa and that instead exclusively studying those open-ended subsystems may always lead to incomplete theories about their origin.
\item Life, and human civilization, owes its existence to the conditions generated by the universe. These include energetically favorable environments, mixing processes, and a complex collection of simple molecules. These conditions did not exist throughout much of the universe's lifetime, although it is unknown exactly which periods could have allowed for the formation of life or other lifelike open-ended processes.
\item Life appears to have arisen quite quickly after the formation of the Earth 4.5 billion years ago, specifically within around 750 million years and possibly hundreds of millions of years earlier based only on known microfossil evidence. This may suggest, in particular if we assume no large differences in historical chance occurrences, that the universe has done the ``hard part" or ``majority" of ab initio biogenesis.
\item The universe was able to take advantage of the (ridiculously) large pool of resources available to eventually generate life and then humans. That is both astonishing, in the sense of defying expectation, and fascinating, in the sense of begging for explanation. Furthermore, there is an observation (or perhaps more aptly, idea) in open-endedness studies\footnote{See Picbreeder.} that open-endedness is the quickest or most efficient process\footnote{With respect to number of ``steps" involved in the process.} to anything. This includes more open-endedness. In other words, this all suggests that an open-ended process like the universe may be the best or most feasible path to life-like open-endedness and, eventually, strong artificial intelligence.\footnote{Of course a version of this exact idea is already in play---human civilization is currently making progress on the AI problem, using open-endedness on the social level.}
\item It may be the case that there are fundamental physical or chemical causes behind life or human civilization that would be difficult to recover or replace in an open-ended system without using related mechanisms.
\item Focusing on ``higher" forms of open-endedness, where abstractions like agency are found, may miss insights more easily or only gained by studying lower forms of open-endedness, where less structure is built into the process.
\item The objects of evolution and human civilization are living organisms\footnote{And arguably viruses.} and humans, respectively. But both open-ended processes involve much more than these entities; in particular the Earth, the Sun, and many other physical influences play a pivotal role that can be difficult to pin down or even lost entirely when focusing solely on these higher-level open-ended processes. It is unclear exactly how much of the local universe should be considered and accounted for when studying these processes. In contrast, the universe itself includes everything with a causal influence, so there are no external environment factors that are not naturally a part of the analysis. In other words, put simply, focusing on the universe places a natural emphasis on environment and context.
\end{itemize}

\section{What drives open-endedness in the universe?}

%%% the universe sets the stage, starts the engine, everything else is "just a consequence", "just follows"
%%% Explain this idea first!!!!

We've identified a number of reasons why it makes sense to first focus on the origin of open-endedness within and at the level of the universe itself, rather than within or at the level of evolution or human civilization. So then why is the universe open-ended? What drives or generates this? Why do endless, compounding, interesting things appear rather than not? Why did these creative processes not peter out at some point? Why did they even start in the first place?

...

% A natural place to start is by considering what proximal causes bring about open-ended phenomena and then trace to the more distal causes.

%%% this may point to a fundamental assumption in labeling the universe as open-ended in the same way as evolution and human cilization; namely, that the open-ended occurs `everywhere', is omnipresent, the factors driving it all-encompassing, and there are no "dead spots", non-open-ended subsets/regions, which should be mentioned earlier
%%% perhaps should first discuss some likely helpful framings first---while Earth is the only region of the universe known to contain open-ended subprocesses, there are reasons to suspect universal factors are at play...in other words, we suspect that it is less useful to ask "why did life arise on Earth? i.e. what's special about Earth?" and more useful to ask "What about the universe brought about the opportunity for life to arise on Earth?"...something like this
%%%% start with the early universe / big bang period? or start with 'most' open-ended phenomena and trace to both other less open-ended phenomena and distal causes
%%% important to keep in mind scale


\subsection{Thermodynamic computing and open-endedness}

Thermodynamic computing is a recently developed idea that (1) hypothesizes that thermodynamics drives the self-organization and evolution of natural systems and (2) proposes that thermodynamics could drive the self-organization and evolution of future computing systems, making them more efficient and capable. It is natural from (1) to further hypothesize that open-endedness in natural systems is simply a consequence of this self-organizing thermodynamic evolution process. Furthermore, perhaps maximally open\footnote{in the sense of thermodynamically open.} thermodynamic computers will be open-ended.

Much more could be said, but the crux of all of this is that thermodynamic computing seems like the most promising avenue to achieving open-endedness. So how do we realize such thermodynamic computers?

%%%%%%

Some thoughts:
\begin{itemize}
\item What is the relationship between thermodynamic computing and open-endedness? Are `maximal' thermodynamic computers necessarily open-ended?
\item What kind of systems maximize entropy production {\em rates} in the environment while minimizing entropy production rates internally? What kind of systems self-organize, adapt, and evolve on useful\footnote{i.e. human} spatial and temporal scales at ordinary temperatures?
\begin{itemize}
    \item Low activation energies, limited potential
    barriers
    \item Certain features held roughly constant; e.g. temperatures, pressures, etc.
\end{itemize}
\item
\end{itemize}

Some additional thoughts towards implementation, especially branching off of existing efforts (in particular the Thermodynamic State Machine Network \cite{...}):

\begin{itemize}
\item Computability is a barrier for the TSMN and indeed all systems implemented on Turing-based computers. In general any thermodynamic computing system implemented on Turing-based computers, which are fundamentally state machines that jump from equilibrium state to equilibrium state by the nature of their operation (and construction), will be merely a simulation, and one limited by the deterministic, stateful nature of a computer.
\item Scalability
\item How connected are equilibration and adaptation in the TSMN? Can the latter effect the former, in a sense? They are connected...
\item TSMN method: define an energy function that changes/updates over time to encourage the finding of ground states and quick transitions between them, then generate the state by sampling the Boltzmann distribution under that energy function, and iterate. Updates to the energy function occur at ground states.
\end{itemize}

\textbf{Is open-endedness merely a product of scale? Is evolution an \textit{ongoing} open-ended process simply because the dynamics are `slow' and the scale is large? To what extent is this the primary difference between these open-ended systems and existing artificial attempts at open-endedness? Then it is about scalability...}

Key insights and messages: Earth needs to be the metaphor and/or model just as much as a living organism. Life, and in particular a living organism, exists in a context.

Concentrating on Earth and the biosphere:
\begin{itemize}
\item Earth is not in a ground state?
\item Earth is not really in connection with a (relevant) thermal reservoir / heat bath. While the external environment, outer space, can absorb energy, it does not supply a significant amount of energy---this comes almost entirely\footnote{99.97\%, see online sources} from the Sun, the source of potential or free energy. Earth's energy and temperature are entirely tied up with its interaction with the Sun. And since the Sun is entirely self-powered, the external source and sink of Earth's energy are approximately disjoint.
\end{itemize}

Model changes:
\begin{itemize}
\item scale, scale, scalability
\item Use dynamics that are valid in a nonequilibrium setting
\item Use more physically-accurate energies, conservation laws, adaptation methods, etc.
\item Continual, recursive, incremental adaptation, evolution, and self-organization should naturally result. Interactions within the system should enable this, with adaptation and organization leading to new types of interaction leading to more adaptation and new structures, etc.
\item More degrees of freedom beyond the transportation and storage of charge?
\item Self-similarity, recursion
\item ``State-energy landscape''---very weak connection in the case of TSMN
\item Organization \textbf{depends} on the transport of conserved quantities, in particular the thermodynamic potentials that drive this transportation.
\item Energy barriers---present in TSMN but inadequate?
\item Both charge and, more importantly, \textbf{energy} are conserved, and energy is ultimately the more important transported quantity.
\item (potential) energy is the criterion/loss function, force is the update gradient (rule), mass and time are source of learning rate
\item Energy is eventually is returned to the external environment but in the meantime it bounces around a system, e.g. Earth, for awhile, while it can still be stored or used for work. \textbf{The lack of energy storage may be a primarily limitation of the TSMN.} Or, if it is storing energy, where does the energy come from? This seems to violate energy conservation, although this is never actually promised. Perhaps not violate conservation, but at least there is a disconnect between the energy-driven dynamics and energy storage, which comes in the form of adaptation when the system enters a ground state. \textbf{This is still unclear---is energy conserved in the TSMN?}
\item \textbf{Energy storage, or specifically free energy storage, is key for open-endedness in nature.} Far-from-equilibrium structures that store free energy during their dissipation activities are an opportunity for other structures to form that take advantage of these child structures.
\item \textbf{Open-endedness emergence}: continual interesting novelty \textleftarrow\, continuous emergence of meaningfully novel (``new") dissipative structures \textleftarrow\, new dissipative structures create new dissipation opportunities, in an continually emerging fashion \textleftarrow\, preexisting dissipative structures are able to adapt to take advantage of these new dissipative opportunities (the emergence of wholly new structures to take advantage of these opportunities becomes increasingly untenable/infeasible as the system and these new structures at the ``cutting-edge" become increasingly sophisticated); highly complex dynamics across varying interacting scales with many (interacting) degrees of freedom + continual input of new free energy to be dissipated
\item In the TSMN there is a partition between the energy function, which constitutes all adaptation features and shapes system dynamics, and system state, specifically the features that make up the system state, namely edges codes/charges. These interact in a sense but are entirely distinct. Put differently and more specifically: transported quantities (which constitute the environment potentials driving change in the system) play no direct role in the adaptation of the system. This stands in direct contrast to natural systems. In such systems energy is the (only?) universal transported quantity and the distribution of energy is directly tied to the structures and mechanisms that constitute and drive adaptation and all other dynamics. For example, almost all energy on Earth arrives via photons, whose number is not conserved but which carry energy that is absorbed as work by a variety of different materials. This changes the state of these materials, changing the global energy landscape, driving further exchanges of energy and changes of state, etc. Throughout this process materials, including charges, are generally conserved but are only transported on highly varying timescales.
\item The TSMN is run in such a way that environment potentials do not operate independently of the TSMN. Each time energy is injected into the system via external input nodes there is a waiting period for the network to find a ground state before any additional energy is input. In this way it is trained much like a neural network or many other standard machine learning methods. There is no such a priori artificial synchronicity or state-dependent forcing in natural systems---e.g. since its formation, Earth is continually hit with incoming energy from the Sun at a roughly constant rate. Its own natural dynamics can introduce state-dependent shielding of parts of the system from external forcings, but this is part of the system itself, part of its dynamics and evolution.
\item The interaction landscape is also very limited in the TSMN! A single binary interaction along each edge. This leads to very limited transportation dynamics.
\item Why a network? Local, finite interactions?
\item Boltzmann dynamics only applies at thermal equilibrium, which is not the correct context. This really does seem like a fundamental mismatch in many ways; think about microstates and macrostates, entropy, the TSMN ostensibly being a system which is driven far from equilibrium, etc. Put differently, all of the far from equilibrium organization is contained in the energy function; the rest, that is, the actual dynamics, are generated from an equilibrium Boltzmann distribution. This at the very least highlights how limited the actual nonequilibrium effects/parts are---there is no interaction except with the equilibrium part. This suggests that the ``self-organization" is highly limited. Does it say more?
\item Thermal equilibrium is necessary to apply the Boltzmann distribution / canonical ensemble. Otherwise it becomes necessary to care about other degrees of freedom besides simply the total energy. In particular given that any drive likely interacts with potential and kinetic energy in different ways.
\item We understand the driven nonequilibrium, driven thermodyanmics, but need to correctly account for the \textbf{work} done by drives/potentials.
\item There seems to be a limited analogy one can setup between the driven-beads experiment and the TSMN---thinking of the codes as charge and the beads, oil, and air as the network.
\item While the potential / drive applied by the Sun is macroscopically appoximately constant and microscopically distributed atemporally, how it is received by the Earth system varies complexly as a function of system state.
\item Important to understand how structure and behavior formation actually happens on Earth (and in the universe): in a nutshell, work from a drive (source of free energy, a potential, a field) raises the \textbf{potential} and/or \textbf{kinetic} energy of constituents of a collection of matter, which then either (1) go on to drive other parts of the system (other matter) or (2) fall back down into lower potential states, releasing energy (often dissipated in the form of heat) and forming new structures, which may play a role in the effects of further drives. Active self-organization in the form of structure formation and maintenance and corresponding ``behavior" is the result of this process continually repeating in a chained way.
\item Locality enforced in the TSMN via the alternating partition update method of the equilibrium algorithm
\item The central relevance of large ground state degeneracy characterizes the TSMN and heavily contrasts it with the Earth system, or more generally the physical world, where global ground states are generally never obtained and therefore are not particularly relevant. Of course that does not mean that degeneracy of higher energy global ground states could not still be relevant.
\item Energy dissipation and potentiation: in the TSMN, energy enters as excitations at ``high potential" nodes or as a result of thermal fluctuations and is dissipated by the transport of an excitation to a ``low potential" node or a complementary excitation. This process should be more closely tied to structure formation. In many real nonequilibrium processes this is key (e.g. biological processes).
\item Earth(-Sun) system is a somewhat extreme / non-archetypal case from the perspective of the thermodynamic computing principles outlined by Hilton, yet its the quintessential example in nature of the core phenomena underlying thermodynamic computing.
\end{itemize}

Normal Computing insights about learning algorithm generalization could be useful. \textbf{Examine their other work for mapping the framework to specific ML algorithms.}

Examine energy conservation in TSMN and consider more generally. Then consider the energy-based framework and whether/how to modify. Also keep in mind reversibility and reversibility. Consider how applied work should be taken into account---how does TSMN do it?

We may be thinking at this point about a different level or scale of ``problem" context than the TSMN is designed to address. Perhaps the physics of the TSMN framework are relatively correct? And it's mainly a question of many missing degrees of freedom, reducing simplifications and approximations, etc. \textbf{So start by briefly laying out the physics of the framework from a super high level.}

Draft model idea:
\begin{itemize}
\item Cylinder lattice 2D $[n] \times [n]$ for $n \in \mathbb{N}$
\item $N$ particles, $k$ different types of counts $n_{1}, \ldots, n_{k}$, $\sum n_{i} = N$, with Pauli exclusion for position within each particle type
\item Particles have masses $m_{1}, \ldots, m_{k}$ and momenta $p_{1}, \ldots, p_{k}$
\item Restrict now to $k = 2$ particle types---this does not actually work as is, likely need more types
\item Short-range linear\footnote{Because we are working in a 2D lattice.} discrete potential, for example:
\begin{align*}
V(\{x_{a_{1}}, y_{a_{1}}, x_{a_{2}}, y_{a_{2}}\}) &= \max\{3-d, 0\} \\
V(\{x_{b_{1}}, y_{b_{1}}, x_{b_{2}}, y_{b_{2}}\}) &= \max\{3-d, 0\} \\
V(\{x_{a}, y_{a}, x_{b}, y_{b}\}) &= \min\{d-3, 0\}
\end{align*}
where $a, a_{1}, a_{2}$ are particles of type 0, $b, b_{1}, b_{2}$ are particles of type 1, $x_{a}$ and $y_{a}$ denote the coordinates of particle $a$, and $d$ is the (lattice) distance between the given particles.
\item The particles are immersed in a viscous heat bath with temperature $T$ and damping coefficient $\lambda$, implying that they obey Langevin dynamics:
\begin{align*}
m\ddot{\mathbf{x}} + \nabla V(x) + \lambda\dot{\mathbf{x}} - \sqrt{2\lambda k_{\textnormal{B}}T}\boldsymbol{\eta}(t) = 0,
\end{align*}
where $m$ is the mass and $R(t)$ is a white noise random variable. We likely assume an overdamped regime where the inertial force term is so much smaller than the remaining terms that we may neglect it. This gives
\begin{align*}
{d\mathbf{x} \over dt} = -{\nabla V(x) \over \lambda} + \sqrt{2 k_{B}T\over \lambda}\boldsymbol{\eta}(t).
\end{align*}
\item For a free given particle, if its kinetic energy KE exceeds its negative potential energy PE then it feels a force as per the Langevin dynamics but continues moving freely. Otherwise, it falls into the potential well and takes on a definite position determined by the potential and Pauli exclusion. The lost kinetic energy is released into the local environment as a uniform or constant wave of spreading energy that travels instantaneously; in other words, all particles within a certain fixed distance receive an impulse proportional to the magnitude of the energy release. Particles adsorb a fraction of this energy corresponding to the fraction of the (discrete) wave front that they occupy (i.e. proportional to their distance from the source), increasing their kinetic energy. The rest is dissipated into the bath. Conversely, for a bound particle that is not confined\footnote{By other particles or an environment boundary.}, if a drive pushes its KE above negative PE then it leaves the well with a velocity corresponding to KE + PE kinetic energy.
%%% What about "molecule" breaking? That is, one bond breakig of a multiply-bonded particle?
%%% For bonded particles how much should this process iterate vs. be transfered into the center of mass of the "molecule"/structure?
\item The external drive works the same way. Energy is injected into the lattice from one side, the closed side, in the form of a wave of spreading energy. Spatially and temporarily the wave has a given random energy distribution, to be determined. Blackbody ``spectrum" could make sense.
\item We account for these energy injections, via the external drive or in-system energy releases, by adding an external force term to the Langevin dynamics,
\begin{align*}
{d\mathbf{x} \over dt} = -{\nabla V(x) \over \lambda} + {F_{ext}(x) \over \lambda} + \sqrt{2 k_{B}T\over \lambda}\boldsymbol{\eta}(t).
\end{align*}
\item Drive/force, potentiation, free energy, energy conservation and transport
\begin{itemize}
    \item
\end{itemize}
\end{itemize}

Draft model idea 2:
\begin{itemize}
\item Coordinate space is a cylindrical lattice $[n] \times [n]$ for $n \in \mathbb{N}$, with a periodic boundary on the first dimension
\item $N$ particles, $k$ different types of counts $n_{1}, \ldots, n_{k}$, $\sum n_{i} = N$, with Pauli exclusion for position within each particle type and across particles $2$ through $k$
\item Particles have masses $m_{1}, \ldots, m_{k}$, charges $-1, 1, \ldots, 1$ denoted by $q_{1}, \ldots, q_{k}$, and momenta $p_{1}, \ldots, p_{k}$
\item Finite-range potential
\begin{align*}
V(\{a, b\}) = \begin{cases}
m_{a}m_{b}q_{a}q_{b}f(d(a, b)), &x \leqslant \rho, \\
0, &\textnormal{else}.
\end{cases}
\end{align*}
where $a$ and $b$ are particles, $d(a, b)$ is the lattice distance between $a$ and $b$, $f$ is some fixed monotonic function, and $\rho$ is a fixed distance specifying the range of the potential (e.g. something like $\rho = 3$). In the setting of electrodynamics or gravitation $f$ would be a logarithm in a 2D space, but simpler forms could also make sense.
\item The particles are immersed in a viscous heat bath with temperature $T$ and damping coefficient $\lambda$, implying that they obey Langevin dynamics:
\begin{align*}
m\ddot{\mathbf{x}} + \nabla V(x) + \lambda\dot{\mathbf{x}} - \sqrt{2\lambda k_{\textnormal{B}}T}\boldsymbol{\eta}(t) = 0,
\end{align*}
where $\mathbf{x}$ is the position vector of a given particle, $m$ is the mass, and $\boldsymbol{\eta}(t)$ is a white noise random variable. We likely assume an overdamped regime where the inertial force term is so much smaller than the remaining terms that we may neglect it. After rearrangement this gives
\begin{align*}
{d\mathbf{x} \over dt} = -{\nabla V(x) \over \lambda} + \sqrt{2 k_{B}T\over \lambda}\boldsymbol{\eta}(t).
\end{align*}
Furthermore, separate from the force generated by the potential, particles may be acted on by external forces or drives, adding an additional term to the Langevin equation:
\begin{align*}
{d\mathbf{x} \over dt} = -{\nabla V(x) \over \lambda} + {F_{ext}(x) \over \lambda} + \sqrt{2 k_{B}T\over \lambda}\boldsymbol{\eta}(t).
\end{align*}
\item Constraints are introduced into the dynamics by Pauli exclusion, which effectively introduce an infinite energy barrier. When a point $x$ is occupied relative to a given particle type $i$ we set
\begin{align*}
V(x) = \infty
\end{align*}
for all particles of type $i$.
\end{itemize}

Draft model idea 3:
\begin{itemize}
\item Coordinate space is a cylindrical lattice $[n] \times [n]$ for $n \in \mathbb{N}$, with a periodic boundary on the first dimension; we write the points as vectors $\mb{r} = (x, y)$
\item $N$ particles, $k$ different types of counts $n_{1}, \ldots, n_{k}$, $\sum n_{i} = N$, with Pauli exclusion for position within each particle type and across particles $2$ through $k$
\item Particles have masses $m_{1}, \ldots, m_{k}$, charges $-1, 1, 2, \ldots, k-1$ denoted by $q_{1}, \ldots, q_{k}$, and momenta $p_{1}, \ldots, p_{k}$
\item Finite-range potential
\begin{align*}
V_{a}(r) = \begin{cases}
q_{a}f(r), &r \leqslant \rho, \\
0, &\textnormal{else}.
\end{cases}
\end{align*}
where $a$ is a particle, $r$ is the lattice distance from $a$, $f$ is some fixed monotonic function, and $\rho$ is a fixed distance specifying the range of the potential (e.g. something like $\rho = 3$). In the setting of electrodynamics or gravitation $f$ would be a logarithm in a 2D space, but simpler forms could also make sense. The potential is additive over multiple sources. The potential energy of a particle $b$ at point $\mb{r}$ is then given by
\begin{align*}
U(\mb{r}; b) = q_{b}V(\mb{r}),
\end{align*}
with the caveat that we do not include any self-interactions, so there is no contribution to the potential energy of $b$ from $b$ itself.
\item Constraints are introduced into the dynamics by Pauli exclusion, which effectively introduce an infinite energy barrier. When a point $\mb{r}$ is occupied relative to a given particle type $i$ we set
\begin{align*}
U(\mb{r}; i) = \infty
\end{align*}
for all particles of type $i$. It is physically natural to determine occupation based on charge polarity; two particles cannot occupy the same position if their charges have the same sign.
\item The particles are immersed in a viscous heat bath with temperature $T$ and collision coefficient $\gamma$.
\item The particles are also subject to a space and time dependent external drive or force field $\mb{F}_{ext}(\mb{r}, t)$. In general this drive could take any form, whether stochastic or deterministic, and interesting behavior might emerge across a variety of drive choices. We choose to focus on physically-inspired drives that evoke the spectrum of solar radiation near Earth. Assuming $\mb{F}_{ext}(\mb{r}, t)$ is generated by blackbody radiation traveling vertically in from the top of the lattice, in accordance with Planck's law we take
\begin{align}
\label{drive_dist}
\mb{F}_{ext}((1, y), t) \sim {F^{3} \over e^{\alpha\beta F} - 1} \quad \textrm{and} \quad \mb{F}_{ext}((x, y), t) = \psi(\mb{r}, t)\mb{F}_{ext}((1, y), t),
\end{align}
where $\alpha$ is a constant tuning the strength of the drive and $\psi(\mb{r}, t)$ specifies how the drive varies in space and time. To emulate solar radiation near Earth, the constant $\alpha$ should be set such that the maximum of the distribution (\ref{drive_dist}) is around an order of magnitude less than the strongest binding energies attained by particles in the system. Radiation cannot in general pass through particles unhindered. Therefore $\psi$ should vary across the lattice in a manner that reflects this screening effect. We could treat particles as perfect barriers or absorbers and set
\begin{align*}
\psi((x, y), t) = \begin{cases}
0, &\textnormal{there exists a particle in } \{x\} \times (y, n], \\
1, &\textnormal{else.}
\end{cases}
\end{align*}
We could also instead consider particles as imperfect barriers and assume an exponential falloff in intensity, setting
\begin{align*}
\psi((x, y), t) = \begin{cases}
\sigma_{1} e^{-\sigma_{2}\ell}, &\textnormal{there exist } \ell \textnormal{ particles in } \{x\} \times (y, n], \\
1, &\textnormal{else.}
\end{cases}
\end{align*}
The constants $\sigma_{1}$ and $\sigma_{2}$ should be set such that an appropriate normalization holds:
\begin{align*}
\sum_{\ell = 0}^{\infty}\sigma_{1}e^{-\sigma_{2}\ell} = 1,
\end{align*}
giving $\sigma_{1} = 1 - e^{-\sigma_{2}}$.
\item  Suppose the particles are in positions $\mb{R}^{(t)} = \{\mb{r}^{(t)}_{1} = (x^{(t)}_{1}, y^{(t)}_{1}), \ldots, \mb{r}^{(t)}_{k} = (x^{(t)}_{k}, y^{(t)}_{k})\}$ with momenta $\Phi^{(t)} = \{p^{(t)}_{1}, \ldots, p^{(t)}_{k}\}$ at time $t$. Since the heat bath is viscous the dynamics are overdamped, meaning that the velocity of any particle almost immediately goes to zero, implying that inertia is not preserved across multiple time steps. So the momenta are completely determined by the current forces acting on the particles: the interaction force, external force, and Brownian fluctuations, all of which at most depend on the positions (when the time is fixed). That is, momentum is not an independent state variable. It thus suffices to select $\mb{R}^{(t+1)}$. The probability of $\mb{R}^{(t+1)}$ is given by the master equation
\begin{align*}
P(\mb{R}^{(t+1)}) = \sum_{\mb{R}}P(\mb{R}^{(t+1)} \mid \mb{R}^{(t)})P(\mb{R}^{(t)} = \mb{R}),
\end{align*}
which reduces to
\begin{align*}
P(\mb{R}^{(t+1)}) = P(\mb{R}^{(t+1)} \mid \mb{R}^{(t)})
\end{align*}
upon fixing $\mb{R}^{(t)}$. So we must determine the conditional distribution $P(\mb{R}^{(t+1)} \mid \mb{R}^{(t)})$.
\item Particles mutually interact with another, so the coordinate variables are not in general independent of one another. This requires us to factorize the conditional distribution into independent pieces. Let $I_{1}, \ldots, I_{\ell}$ be a pairwise independent partition of the particles; that is, such that no particle in $I_{i}$ interacts with any other particle in $I_{i}$ for all $i$. Then we may use Gibbs sampling over the partition to determine $P(\mb{R}^{(t+1)} \mid \mb{R}^{(t)})$. When sampling the appropriate conditional distribution of an independent subset $I_{i}$ of particles we can sample the next state of each particle in $I_{i}$ independently.
\item This reduces the problem to sampling the next state of a single particle. Let $\mb{R}$ be the current state (during the global sampling procedure) with particle $i$ at $\mb{r}^{(t)}_{i} = (x^{(t)}_{i}, y^{(t)}_{i})$ and $\mb{R}'$ be the state with particle $i$ at its next point $\mb{r}^{(t+1)}_{i} = (x^{(t+1)}_{i}, y^{(t+1)}_{i})$ and all other particles fixed. Then we seek an estimate for $P(\mb{R}' \mid \mb{R})$. Two factors govern this distribution. First, particles travel with finite velocity and, according to overdamped Langevin dynamics, their velocity is determined by the forces currently acting on them, including the ``force" produced by the interaction potential $V$. In other words, the net force $f_{net}$ on a particle at time $t$ defines how far it can travel in the single time step from $t$ to $t+1$:
\begin{align*}
d = {f_{net}\Delta t \over \gamma m_{i}},
\end{align*}
where $\Delta t$ is a unit defining the conversion between velocity and a single time step. This bound is not actually correct because the external force field acts instantaneously as a particles moves, so its velocity is path-dependent, but for the purposes of simplicity and efficiency we use this approximation. If $\mb{R}'$ and $\mb{R}$ defer by a particle moving more than $d$ distance then
\begin{align*}
P(\mb{R}' \mid \mb{R}) = 0.
\end{align*}
Now suppose $\mb{R}'$ satisfies this velocity bound. When the particle moves it may encounter energetic influences effecting its probability of reaching a distant point. We imagine that it is taking a single shortest path on its way and track the energies encountered along the way to determine their effect on its probability of reaching the point. Let $\mc{P} = (\mb{R}_{1}, \ldots, \mb{R}_{m})$ be a fixed shortest path between $\mb{R} = \mb{R}_{1}$ and $\mb{R}' = \mb{R}_{m}$ in state space. Following non-equilibrium Boltzmann thermodynamics we take
\begin{align}
\label{transition_prob_eq}
P(\mb{R}' \mid \mb{R}) = {1 \over Z}e^{\beta(W_{\mc{P}} - \Delta E_{\mb{R}\mb{R}'} - B_{\mc{P}})},
\end{align}
where $\Delta E_{\mb{R}\mb{R}'} = U(\mb{R}') - U(\mb{R})$ is the energy difference between $\mb{R}$ and $\mb{R}'$,
\begin{align*}
B_{\mc{P}} = \max\{U(\mb{R}_{i})\}_{i=1}^{m} - \max\{U(\mb{R}), U(\mb{R}')\}
\end{align*}
is the height of the energy barrier encountered while traversing $\mc{P}$, $W_{\mc{P}}$ is the work absorbed during the traversal of $\mc{P}$, $\beta = 1/k_{\B}T$, and $Z$ is the partition function
\begin{align*}
Z = \sum_{\mb{R}'}e^{\beta(W_{\mc{P}} - \Delta E_{\mb{R}\mb{R}'} - B_{\mc{P}})}.
\end{align*}
The absorbed work depends both on the external drive and the Brownian dynamics induced by the bath:
\begin{align*}
W_{\mc{P}} = W^{ext}_{\mc{P}} + W^{rand}_{\mc{P}}.
\end{align*}
The former is given by
\begin{align*}
W^{ext}_{\mc{P}} = \sum_{j=1}^{m-1}\mb{F}_{ext}(\mb{r}_{j}, t) \cdot \Delta\mathbf{r}_{j}
\end{align*}
with $\mb{r}_{1}, \ldots, \mb{r}_{m}$ the sequence of points the particle traverses on path $\mc{P}$ and $\Delta \mb{r}_{j} = \mb{r}_{j+1} - \mb{r}_{j}$. Note that we are assuming the external force acts instantaneously, meaning the moving particle interacts with the force throughout its motion across a single time step. In contrast, we do not assume the Brownian force acts instantaneously, but rather that it is only imparted at the original position of the particle. By Langevin dynamics this force is
\begin{align*}
\mb{f}_{rand} \sim \sqrt{2\gamma m_{i} \over \beta}\boldsymbol{\eta}(t),
\end{align*}
where $\boldsymbol{\eta}(t)$ is a stationary Gaussian process satisfying
\begin{align*}
\langle \boldsymbol{\eta}(t) \rangle &= 0 \\
\langle \boldsymbol{\eta}(t)\boldsymbol{\eta}(t') \rangle &= \delta(t - t')
\end{align*}
for $\delta$ the Dirac delta function. Then
\begin{align*}
W^{rand}_{\mc{P}} = \mb{f}_{rand} \cdot \Delta\mb{r},
\end{align*}
where $\Delta\mb{r}$ is the displacement from $\mb{r}_{i}^{(t)}$ to $\mb{r}_{i}^{(t+1)}$. Summing, we have
\begin{align*}
W_{\mc{P}} = W^{rand}_{\mc{P}} + W^{ext}_{\mc{P}} = \mb{f}_{rand} \cdot \Delta\mb{r} + \sum_{j=1}^{m-1}\mb{F}_{ext}(\mb{r}_{j}, t) \cdot \Delta\mb{r}_{j}.
\end{align*}
\item For a continuous (scalar) potential $P$ the (unit) force or field due to the potential is given by the negative gradient of $P$, $-\nabla P$. The state space considered here is discrete and so the potential $V$ is as well, so we require a different route to defining the force $\mb{F}_{int}$ produced by $V$. There is a natural choice that fits with our expression (\ref{transition_prob_eq}) for the transition dynamics, namely using finite differences of $V$. Write
\begin{align*}
\Delta_{x}V(x, y) = V(x+1, y) - V(x, y)
\end{align*}
for the forward difference operator with respect to $x$ and
\begin{align*}
\nabla_{x}V(x, y) = V(x, y) - V(x-1, y)
\end{align*}
for the backward difference operator with respect to $x$, respecting periodicity. Define similar operators with respect to $y$:
\begin{align*}
\Delta_{y}V(x, y) = \begin{cases}
0, &y = n, \\
V(x, y+1) - V(x, y),  &\textnormal{else,}
\end{cases}
\end{align*}
for the forward difference operator with respect to $x$ and
\begin{align*}
\nabla_{y}V(x, y) = \begin{cases}
0, &y = 1, \\
V(x, y) - V(x, y-1),  &\textnormal{else.}
\end{cases}
\end{align*}
These operators combine to give a notion of gradient for a discrete vector field such as $V$,
\begin{align*}
\nabla V(\mb{r}) &= \left({\Delta_{x}V(\mb{r}) + \nabla_{x}V(\mb{r}) \over 2}, {\Delta_{y}V(\mb{r}) + \nabla_{y}V(\mb{r}) \over 2}\right),
\end{align*}
which we extend linearly to the potential energy $U$. In analogy with the continuous case we define $\mb{F}_{int} = -\nabla U$ to be the force field corresponding to the potential $V$. This is used to compute the net force experienced by particles and thereby determine the velocity bound during a transition.
% only four adjacent lattice sites? or also diagonal ones as well?
%%% explain relation to energy difference in transition probability
\item Fixing a threshold $\upsilon \leqslant 0$, we say a particle $a$ at point $\mb{r}$ is {\em in a bound state} if $U(\mb{r}; a) < \upsilon$.
\item The above gives the absorbed work, but some work is also not absorbed. The total work done by the forces during the time step is
\begin{align*}
W_{total} = f_{rand}d + F_{ext}(\mb{r}_{i}^{(t+1)}, t)(d - \Delta r) + \sum_{j=1}^{m-1}F_{ext}(\mb{r}_{j}, t) \Delta r_{j},
\end{align*}
since $\mb{r}_{m} = \mb{r}_{i}^{(t+1)}$, leaving
%%% where $v = |\mb{v}|$ for any vector $\mb{v}$
\begin{align*}
W_{ex} &= W_{total} - W_{\mc{P}} \\
&= \mb{f}_{rand} \cdot (d - \Delta\mb{r}) + F_{ext}(\mb{r}_{i}^{(t+1)}, t)(d - \Delta r) + \sum_{j=1}^{m-1}\mb{F}_{ext}(\mb{r}_{j}, t) \cdot (\Delta r_{j} - \Delta\mb{r}_{j})
\end{align*}
work unaccounted for. It is clear that, in expectation, the only reason why we would have $W_{ex} > 0$ is because the particle is in a bound state in $\mb{R}'$. Accordingly, in the event the system is selected to transition to $\mb{R}'$, this energy could either be transferred through the bond into the bound state or it could be emitted out into the surrounding environment; in general we would expect some combination of these two possibilities. In real physical systems the latter generally occurs due to photon or phonon emission events, which require energies near the energy scale of the bond\footnote{Check.}. We use this criterion to determine how the extra work is transformed. First, we set a threshold $\mu$ specifying the minimum energy required to produce an effect; if any amount $W$ of energy satisfies
\begin{align*}
W \leqslant \mu
\end{align*}
then that energy is discarded from the dynamics; we imagine that it is absorbed by the heat bath without measurably effecting any particles. This threshold $\mu$ could be e.g. set based on the minimum amount of energy required to, in expectation, move a particle one unit,\footnote{Or produce some other minimum change in the system.} and could be fixed or vary depending on the context. Now, if %% maybe also for small work?
\begin{align*}
W_{ex} \geqslant \epsilon U(\mb{r}^{(t+1)}_{i}; i)
\end{align*}
for some fixed constant $\epsilon \approx 1$, the work is released as a uniform finite range energy wave. Each particle $a$ within a distance $\delta_{1}$ receives kinetic energy
\begin{align*}
K_{1} = {\epsilon_{1}W_{ex} \over |\mb{r} - \mb{r}_{i}^{(t+1)}|},
\end{align*}
where $\mb{r}$ is the position of $a$ and $\epsilon_{1} > 0$ is a fixed constant specifying the fraction of the wave's energy the particle receives and satisfying % kinetic energy, momentum?
\begin{align*}
\sum_{|\mb{r} - \mb{r}_{i}^{(t+1)}| \leqslant \delta_{1}}{1 \over |\mb{r} - \mb{r}_{i}^{(t+1)}|} = {1 \over \epsilon_{1}}
\end{align*}
for energy conservation, where the sum is over all points in the $\delta_{1}$-ball, making $\epsilon_{1}$ a simple function of $\delta_{1}$. The range parameter $\delta_{1}$ could be a small constant, reflecting fast absorption of traveling energy by the bath, or it could be set so that at the boundary the energy at a point drops below the minimum energy threshold $\mu$. We treat the wave as being released after the end of the current timestep, so the particles which receive this energy are determined after all state updates at time $t$, including both particle and bound state updates\footnote{Meaning emissions only direct effect particles.}. There are many ways a particle could receive this kinetic energy; we choose perhaps the simplest setup that explicitly conserves momentum. At its next state update, we add
\begin{align*}
\max\{K_{1}(\mb{n} \cdot \mb{s}), K_{1}\}
\end{align*}
to the absorbed work $W_{\mc{P}}$ of particle $a$, where $\mb{n}$ is a unit vector in the direction of $\mb{r} - \mb{r}_{i}^{(t+1)}$ and $\mb{s}$ is a vector corresponding to the proposed state update on $a$, and add
\begin{align*}
\max\{K_{1} - K_{1}(\mb{n} \cdot \mb{s}), 0\}
\end{align*}
to the corresponding excess work $W_{ex}$. It is also necessary to update the velocity bound; we add
\begin{align*}
\Delta t\sqrt{2K_{1} \over m_{a}}
\end{align*}
to the distance the particle may travel.
%%% Should convert this to include particle masking, just like with the drive
\item It remains to handle the case where
\begin{align*}
\mu \leqslant W_{ex} < \epsilon U(\mb{r}_{i}^{(t+1)}; i)
\end{align*}
and the energy is transferred into the bound state $A$ of particle $i$. To do this we add a second phase to our update algorithm which employs blocked Gibbs sampling over bound states. We run this phase after each phase of particle-level state sampling, meaning that the time has already incremented by 1 and we are simply refining the position samples at the current time step.\footnote{Could consider some mixed strategy, if needed. Update: in fact this is needed, see a much later bullet point.} For the purposes of efficient sampling we treat particles in a bound state as perfectly correlated, fixing their positions relative to one another. This reduces the state variables to the position and momentum of the center of mass of a bound state $A$, as well as the orientation and angular momentum about its center of mass. Since space is a square lattice the allowed orientations are $n\pi/2$ radian rotations of the current orientation of $A$, where $n$ is an integer.
\item Particle $a$ at position $\mb{r}$ is {\em bonded} to a nonempty particle set $B$ during an update at time $t+1$ if $a$ is in a bound state,
\begin{align*}
-q_{a}\sum_{b \in B}V_{b}(\mb{r}) \geqslant W_{ex} + \omega  %% extend this
\end{align*}
for some $\omega \geqslant 0$, and this inequality is false for all proper subsets of $B$, that is, $B$ is minimal. Furthermore,  $a$ can only be bonded to mutually disjoint sets. When this relation is extended to an equivalence relation, we call the equivalence classes {\em bound states}.\footnote{Analogous to molecules.}
%% what about a few weaker "bonds" overpowering a single strong bond? Need to think about when exactly there is a bond---what exactly is the maximal collection of particles defining a single bound state; is adjacency (i.e. being fully in a potential well) required? Probably not...
%% could consider making it also relative to excess work level
%% Particles can be bonded to multiple other particles!
%% It should be relative to (excess) work level, because we are not actually choosing what is bonded / bound in some absolute sense, we are choosing what would *stay bonded/bound* in a given time step, and this is clearly dependent on the applied work
%% Consider making omega stochastic
%% Transfer to other particles!
\item We treat the heat bath as sufficiently viscous to maintain overdamped dynamics even at the larger scale of blocked states, so the momenta and angular momenta are both completely determined by the current forces acting on the particles. These forces arise from the aggregate excess work $W_{ex}$ of all the particles in the bound state. Since angular momentum is a factor it is necessary to track the directions of the forces, meaning that we must return to analyzing in detail the forces on the individual particles.

The force applied on the bound state is determined by the force experienced by a particle when it comes to rest in the bound state. Write $\mb{f}^{(t)}_{i}$ for this force, the last force experienced by particle $i$ at time $t$, determined by summing the Brownian force and final external force:
%%% and any other forces
\begin{align*}
\mb{f}^{(t)}_{i} = \mb{f}_{rand} + \mb{F}_{ext}(\mb{r}_{i}^{(t+1)}, t).
\end{align*}
Note that, unlike previously when calculating the velocity bound, we are determining this net force {\em after sampling}, so this is a single efficient computation. Writing $m_{A} = \sum_{i \in A}m_{i}$ for the mass of $A$, $\mb{r}_{A}$ for the position of the center of mass of $A$, $\mb{f}_{A}$ for the force on $A$, and $\boldsymbol{\tau}_{A}$ for the torque on $A$, we have
\begin{align*}
\mb{f}_{A} &= \sum_{i \in A}\mb{f}_{i}^{(t)} \\
\boldsymbol{\tau}_{A} &= \sum_{i \in A}\boldsymbol{\tau}_{A,i} := \sum_{i \in A}(\mb{r}_{i}^{(t+1)} - \mb{r}_{A}) \times \mb{f}_{i}^{(t)}.
\end{align*}
Since orientation is a degree of freedom we have both a linear velocity bound and angular velocity bound on the possible next states of $A$. These bounds are not only determined by the forces exerted on $A$ but also on the amount of work those forces have already expended moving the constituent particles of $A$. We begin by stating the bounds obtained without this second factor:
\begin{align*}
d = {f_{net}\Delta t \over \gamma m_{A}} \quad \textrm{and} \quad \theta = {\tau_{net}\Delta t \over \gamma m_{A}}.
\end{align*}
Letting $\mb{R}$ be the current state with $A$ at position $\mb{r}_{A}$ and orientation $\theta$\footnote{Without loss of generality we may assume $\theta = 0$.} and $\mb{R}'$ be the state with $A$ at its next point $\mb{r}_{A}'$ and orientation $\theta'$ and all other bound states fixed, as before if $\mb{R}'$ and $\mb{R}$ defer by a bound state moving more than $d$ distance or rotating more than $\theta$ radians then
\begin{align*}
P(\mb{R}' \mid \mb{R}) = 0.
\end{align*}
Assume $\mb{R}'$ satisfies these bounds and let $\mc{P} = (\mb{R}_{1}, \ldots, \mb{R}_{m})$ be a fixed shortest path between $\mb{R} = \mb{R}_{1}$ and $\mb{R}' = \mb{R}_{m}$ in state space. Then we take
\begin{align*}
P(\mb{R}' \mid \mb{R}) =
{1 \over Z}e^{\beta(\max\{W_{\mc{P}}, W\} - \Delta E_{\mb{R}\mb{R}'} - B_{\mc{P}})},
\end{align*}
where
\begin{align*}
W = \sum_{i \in A}W_{i, ex}
\end{align*}
is the sum of the extra work terms corresponding to each particle of $A$. The work $W_{\mc{P}}$ absorbed in the traversal of $\mc{P}$ is
\begin{align*}
W_{\mc{P}} = \mb{f}_{A} \cdot \Delta \mb{r}_{A} + \sum_{j=1}^{m-1}\boldsymbol{\tau}_{A}\Delta \theta_{j} = \mb{f}_{A} \cdot \Delta \mb{r}_{A} + \boldsymbol{\tau}_{A}\Delta \theta,
\end{align*}
with $\theta_{1}, \ldots, \theta_{m}$ the sequence of orientations $A$ attains on the points of path $\mc{P}$, $\Delta \theta_{j} = \theta_{j+1} - \theta_{j}$, $\Delta\theta$ the total angular displacement from $\theta$ to $\theta'$, and $\Delta\mb{r}_{A} = \mb{r}_{A'} - \mb{r}_{A}$. Any excess work $W - W_{\mc{P}}$ remaining after the update is transferred to any newly bonded particles in proportion to the strength of the bond. Suppose the update to $\mb{R}'$ is selected. Let $B$ be the set of all particles not in $A$ that are bonded to particles in $A$. By definition these particles were not bonded to any particles in $A$ in state $\mb{R}$. Let $b \in B$ and define
\begin{align*}
U_{b}(A) = \sum_{a \in A}q_{a}V_{b}(\mb{r}_{a}),
\end{align*}
the cumulative potential energy of $A$ due to $b$. Writing $\Delta U_{b}(A)$ for the difference in this quantity between $\mb{R}$ and $\mb{R}'$, particle $b \in B$ receives kinetic energy
\begin{align*}
K_{2} = {\Delta U_{b}(A) \over \sum_{b \in B} \Delta U_{b}(A)}(W - W_{\mc{P}}).
\end{align*}
We could simplify the energy transfer by distributing this among all particles within interaction range, rather than only those that bond with $A$ or that are in the larger bound state of $A$ in $\mb{R}'$. The absorption of this energy works the same way as with the energy waves described earlier. At its next state update, we add
\begin{align*}
\max\{K_{2}(\mb{n} \cdot \mb{s}), K_{2}\}
\end{align*}
to the absorbed work $W_{\mc{P}}$ of particle $b$, where $\mb{n}$ is a unit vector in the direction of $\mb{f}_{A} + \mb{f}_{torque}$ with
\begin{align*}
\mb{f}_{torque} = {\mb{r} - \mb{r}_{A} \over |\mb{r} - \mb{r}_{A}|^{2}} \times \boldsymbol{\tau}_{A},
\end{align*}
the force applied on $b$ by the torque, and $\mb{s}$ is a vector corresponding to the proposed state update on $a$. As previously we add the remaining kinetic energy to the corresponding excess work $W_{ex}$ and update the velocity bound appropriately.
\item There is a complication that we need to address. Due to the need for Gibbs sampling not all particles are sampled simultaneously. It is therefore possible that the bound state a particle is in when its next state is sampled may be different from the bound state it is in at the end of the particle-level update phase. This means that we cannot assume work is transferred from a particle to this later bound state in the simple manner suggested above. There are several ways we could deal with this, including interleaving block and particle-level Gibbs updates. We choose a more efficient route that avoids additional sampling while maintaining consistent physics. We treat the work as first transferred directly to each particle. After all particle updates any work remaining is then applied to the bound state. Consider a particle $a \neq i$ in the bound state $A$ of $i$ immediately after the state of $i$ is updated. We must determine the force $\mb{f}_{a}$ on $a$ and work $W_{a,i}$ transferred to $a$ due to the excess work of particle $i$. By Langevin dynamics we have
\begin{align*}
\mb{f}_{a} = \mb{f}^{(t)}_{i} + m_{A}\mb{I}_{A}^{-1}(\boldsymbol{\tau}_{A, i} \times (\mb{r}_{a} - \mb{r}_{i})),
\end{align*}
where $\mb{I}_{A}$ is the moment of inertia of $A$. The excess work $W_{a,i}$ transferred to $a$ is then determined by this force in combination with the appropriate velocity bound, which is determined by the velocity bound for the bound state $A$:
\begin{align*}
d = {f_{net}\Delta t \over \gamma m_{A}}.
\end{align*}
\item It remains to determine how to sample from the conditional distribution $P(\mb{R}' \mid \mb{R})$ in either phase of the algorithm. Given the finite, and potentially quite small, range of the potentials and velocities involved, it may be feasible to directly compute the partition function $Z$ or some approximation to it, enabling direct sampling from $P(\mb{R}' \mid \mb{R})$. Failing that, there are several alternatives. We assume it is only feasible to compute the unnormalized density $w(\mb{R}' \mid \mb{R}) = ZP(\mb{R}' \mid \mb{R})$ for each $\mb{R}'$. We begin by factorizing the distribution into a proposal distribution $g$ and an acceptance distribution $A$,
\begin{align*}
P(\mb{R}' \mid \mb{R}) = g(\mb{R}' \mid \mb{R})A(\mb{R}', \mb{R}).
\end{align*}
A natural choice for $g$ would be a finite uniform distribution centered around $\mb{R}$ with range given by the velocity bound. For the acceptance distribution $A$, we could consider simply truncating the density,
\begin{align*}
A(\mb{R}', \mb{R}) = \min\{1, w(\mb{R}' \mid \mb{R})\}.
\end{align*}
While this changes the distribution, there are some properties that are preserved: (1) states with equal probability under $P$ still have equal probability under $A$ and (2) as you approach a ``local" maximum of $P$, the truncated distribution converges to the true distribution, suggesting that this may also be true everywhere else in the long time limit. This form of an acceptance distribution also benefits from values of $w$ commonly lying below 1, which we certainly expect for the Boltzmannian densities we use. We can decrease the impact of the truncation by multiplying by a constant factor $\kappa \leqslant 1$:
\begin{align*}
A(\mb{R}', \mb{R}) = \min\{1, \kappa w(\mb{R}' \mid \mb{R})\}.
\end{align*}
In the limit as $\kappa$ approaches the reciprocal of a least upper bound for the density, $M = \max_{\mb{R}'}w(\mb{R}' \mid \mb{R})$, this method converges to rejection sampling, which does produce the correct distribution:
\begin{align*}
A(\mb{R}', \mb{R}) = \min\left\{1, {w(\mb{R}' \mid \mb{R}) \over M}\right\} = {w(\mb{R}' \mid \mb{R}) \over M}.
\end{align*}

%%% remaining:
%%% - define MCMC algorithm (?)
%%% - Does the drive reflect off the "ground" / lower boundary?
%%% - What about particles at the boundary?

%%% implementation notes
%%% - take into account previous bound states to save unnecessary computations; particles in the same position or with 0 excess work (free particles) are assumed to be in the same bound state
%%% -

%% check energy conservation, lost work? etc.
%%% physical accuracy would have Brownian work only go into movement, while external drive can sent out energy "waves" or "bits"
%%% under this setup only beta-scale weak "bonds" can break between two non-singular groups of particles; otherwise particles break from bonds individually or not at all.
%%% define the stochastic dynamics scheme, MCMC, subset of independent particles at a time; we select a candidate next state using drive-modified (Brownian-randomized?) Boltzman transition probabilities/rates; then apply discrete Langevin dynamics for acceptance method
\end{itemize}


Draft model idea 4:
\begin{itemize}
\item Coordinate space is a cylindrical lattice $[n] \times [n]$ for $n \in \mathbb{N}$, with a periodic boundary on the first dimension; we write the points as vectors $\mb{r} = (x, y)$
\item The space is populated by $N$ particles of $k$ different types of counts $n_{1}, \ldots, n_{k}$, $\sum n_{i} = N$, with Pauli exclusion amongst all particles
\item Particle types $1, \ldots, k$ have masses $m_{1}, \ldots, m_{k}$ and charges $-1, 1, 2, \ldots, k-1$, denoted by $q_{1}, \ldots, q_{k}$. Each particle $a$ has a position $\mb{r}_{a}$ and momentum $\mb{p}_{a}$.
\item Particles produce a finite-range potential
\begin{align*}
V_{a}(\mb{r}) = \begin{cases}
q_{a}, &|\mb{r} - \mb{r}_{a}| = 1, \\
0, &\textnormal{else},
\end{cases}
\end{align*}
where $a$ is a particle and $|\mb{r}| = |x| + |y|$ is the 1-norm of $\mb{r} = (x, y)$. Setting $V_{a}(\mb{0}) = 0$ in this definition naturally excludes self-interactions. The potential is additive over multiple sources, that is,
\begin{align*}
V(\mb{r}) = \sum_{a}V_{a}(\mb{r}).
\end{align*}
The potential energy of a particle $b$ at point $\mb{r}$ is then given by
\begin{align*}
U(\mb{r}; b) = q_{b}V(\mb{r}).
\end{align*}
Pauli exclusion could be expressed by setting $U(\mb{r}; b) = \infty$ for all particles in a multiply-occupied position. To summarize, particles interact with their four nearest neighbor positions and the interaction strength is additive over neighboring particles and multiplicative with their charges.
%%% define internal potential / interaction energy of a set of particles
%%% define (total) energy of a set of particles (PE + internal energy, namely kinetic)
%%% bound state: energy is negative
\item The kinetic energy $K(a)$ of a particle $a$ is given by $|\mb{p}_{a}|^{2}/2m$. The energy $E(A)$ of a set $A = \{a_{1}, \ldots, a_{\ell}\}$ of particles is the sum of the kinetic and potential energies,
\begin{align*}
E(A) = K(A) + U(A).
\end{align*}
The kinetic energy is additive over particles, while the potential energy is given by
\begin{align*}
U(A) &= \sum_{a \in A}\sum_{b}q_{a}V_{b}(\mb{r}_{a}) \\
&= \sum_{a \in A}U(a) - \sum_{i < j}q_{i}V_{j}(\mb{r}_{i}) \\
&= \sum_{a \in A}U(a) - {1 \over 2}\sum_{a \in A}U_{A}(a) \\
&= \sum_{a \in A}U(a) - {U_{A}(a) \over 2},
\end{align*}
where $i$ and $j$ refer to $a_{i}$ and $a_{j}$,
\begin{align*}
U_{A}(a) = \sum_{b \in A}q_{a}V_{b}(\mb{r}_{a}),
\end{align*}
and we may write $U(a) = U(\mb{r}_{a}; a)$ because the positions are fixed. The expression
\begin{align*}
E_{int}(A) := {1 \over 2}\sum_{a \in A}U_{A}(a)
\end{align*}
gives the interaction energy of the particles in $A$. We will ultimately not need these full derivations because the dynamics depend only on changes in energy.
\item We say that a particle $a$ at point $\mb{r}$ is {\em in a bound state} or simply {\em bound} if $E(a) < 0$ or, equivalently, $-U(\mb{r}; a) > K(a)$. In this case we say that particle $a$ is {\em bonded} to each of its bound neighbors of opposite charge sign, that is, every particle $b$ such that
\begin{align*}
U_{b}(\mb{r}; a) := q_{a}V_{b}(\mb{r}) < 0.
\end{align*}
When this relation is extended to an equivalence relation, we call the equivalence classes {\em bound states}.\footnote{Analogous to molecules.}
\item The particles are immersed in a viscous heat bath with temperature $T = 1/\beta$ and collision coefficient $\gamma$. Viscosity introduces both damping and fluctuation, with the latter appearing as a Brownian force
\begin{align*}
\mb{f}_{rand} \sim \sqrt{2\gamma m \over \beta}\boldsymbol{\eta},
\end{align*}
where $\boldsymbol{\eta}$ is a 2-dimensional standard normal distribution and $m$ is the mass of a particle subject to the force.
\item For a continuous scalar field $P$ of potential energies the force due to $P$ is given by the negative gradient of $P$, $-\nabla P$. The state space considered here is discrete, so we require a different route to defining the force produced by the potential energy $U$. We first note that we require an expression for the interaction force on a particle $a$ in order to determine the change in momentum under a particular change $\mb{e}$ in position. To that end it is natural to define a directional gradient $\nabla_{\mb{e}} U(\mb{r}; a)$ of the potential energy of $a$ as the vector satisfying the equation
\begin{align*}
U(\mb{r} + \mb{e}; a) - U(\mb{r}; a) = \nabla_{\mb{e}} U(\mb{r}; a) \cdot \mb{e}.
\end{align*}
This mirrors one standard definition of the gradient via the total derivative. For unit shift vectors $\mb{e}$, we have
\begin{align*}
\nabla_{\mb{e}} U(\mb{r}; a) = (U(\mb{r} + \mb{e}; a) - U(\mb{r}; a))\mb{e}.
\end{align*}
\item The particles are subject to up to two space-dependent external drives represented by fields $\mb{F}_{charge}$ and $\mb{F}_{mass}$ coupled to particle charge and mass, respectively. That is, the force generated by each field on a particle of charge $q$ and mass $m$ is $q\mb{F}_{charge}$ and $m\mb{F}_{mass}$, respectively. In general these drives could take any form, whether stochastic or deterministic, and interesting behavior might emerge across a variety of drive choices. We choose to focus on physically-inspired drives that evoke the spectrum of solar radiation near Earth as well as its gravitational field, capturing the dominant energy sources on Earth. The charge-coupled field is modeled after electromagnetism. We write
\begin{align*}
\mb{F}_{charge} = \mb{E} + {\mb{p} \over \kappa m} \times \mb{B},
\end{align*}
where $\mb{p}$ is the momentum of a particle influenced by $\mb{F}$ and $\kappa$ is a constant determining the relative strength of $\mb{B}$ and $\mb{E}$. Assuming $\mb{E}$ and $\mb{B}$ are generated by blackbody radiation traveling vertically in from the top of the lattice, in accordance with Planck's law we take
\begin{align*}
\mb{E}((1, y)) &= (\theta, 0), \\
\theta &\sim {x^{3} \over e^{\alpha x} - 1}, \textrm{ and} \\
\mb{F}_{ext}((x, y)) &= \psi((x, y))\mb{F}_{ext}((1, y)),
\end{align*}
where $\alpha$ is a constant tuning the strength of the drive and $\psi$ specifies how the drive varies in space. We define $\mb{B}$ as the field orthogonal to both coordinate dimensions of the lattice and of equal magnitude to $\mb{E}$ such that $\mb{E} \times \mb{B}$ points directly downwards in the vertical axis. To emulate solar radiation near Earth, the constant $\alpha$ should be set such that maximum of the distribution of $\theta$ is around an order of magnitude less than the strongest binding energies attained by particles in the system. Radiation cannot in general pass through particles unhindered. Therefore $\psi$ should vary across the lattice in a manner that reflects this screening effect. We could treat particles as perfect barriers or absorbers and set
\begin{align*}
\psi((x, y)) = \begin{cases}
0, &\textnormal{there exists a particle in } \{x\} \times (y, n], \\
1, &\textnormal{else}.
\end{cases}
\end{align*}
We could also instead consider particles as imperfect barriers and assume an exponential falloff in intensity, setting
\begin{align*}
\psi((x, y)) = \begin{cases}
\sigma_{1}e^{-\sigma_{2}\ell}, &\textnormal{there exist } \ell > 0 \textnormal{ particles in } \{x\} \times (y, n], \\
1, &\textnormal{else}.
\end{cases}
\end{align*}
The constants $\sigma_{1}$ and $\sigma_{2}$ should be set such that an appropriate normalization holds:
\begin{align*}
\sum_{\ell = 0}^{\infty} \sigma_{1}e^{-\sigma_{2}\ell} = 1,
\end{align*}
giving $\sigma_{1} = 1 - e^{-\sigma_{2}}$. The mass-coupled field $\mb{F}_{mass}$ is a uniform field modeled after gravitation. Fixing a constant $g$, we set
\begin{align*}
\mb{F}_{mass}(\mb{r}) = g
\end{align*}
everywhere. The constant $g$ should be set such that in general $q\mb{F}_{charge}$ dwarfs $m\mb{F}_{mass}$ unless $m$ is much larger than $q$, although one could also consider alternative schemes.
\begin{comment}
\item The particles are subject to a space-dependent external drive represented by a field $\mb{F}_{ext}$ which is coupled to particle charge; that is, the force generated by the field on a particle of charge $q$ is $q\mb{F}_{ext}$. In general this drive could take any form, whether stochastic or deterministic, and interesting behavior might emerge across a variety of drive choices. We choose to focus on physically-inspired drives that evoke the spectrum of solar radiation near Earth. Assuming $\mb{F}_{ext}$ is generated by blackbody radiation traveling vertically in from the top of the lattice, in accordance with Planck's law we take
\begin{align}
\label{drive_eq}
\mb{F}_{ext}((1, y)) \sim {F^{3} \over e^{\alpha F} - 1} \quad \textnormal{and} \quad \mb{F}_{ext}((x, y)) = \psi((x, y))\mb{F}_{ext}((1, y)),
\end{align}
where $\alpha$ is a constant tuning the strength of the drive and $\psi$ specifies how the drive varies in space. To emulate solar radiation near Earth, the constant $\alpha$ should be set such that maximum of the distribution (\ref{drive_eq}) is around an order of magnitude less than the strongest binding energies attained by particles in the system. Radiation cannot in general pass through particles unhindered. Therefore $\psi$ should vary across the lattice in a manner that reflects this screening effect. We could treat particles as perfect barriers or absorbers and set
\begin{align*}
\psi((x, y)) = \begin{cases}
0, &\textnormal{there exists a particle in } \{x\} \times (y, n], \\
1, &\textnormal{else}.
\end{cases}
\end{align*}
We could also instead consider particles as imperfect barriers and assume an exponential falloff in intensity, setting
\begin{align*}
\psi((x, y)) = \begin{cases}
\sigma_{1}e^{-\sigma_{2}\ell}, &\textnormal{there exist } \ell > 0 \textnormal{ particles in } \{x\} \times (y, n], \\
1, &\textnormal{else}.
\end{cases}
\end{align*}
The constants $\sigma_{1}$ and $\sigma_{2}$ should be set such that an appropriate normalization holds:
\begin{align*}
\sum_{\ell = 0}^{\infty} \sigma_{1}e^{-\sigma_{2}\ell} = 1,
\end{align*}
giving $\sigma_{1} = 1 - e^{-\sigma_{2}}$.
\end{comment}
\item Energy from external drives may flow through the system in several different ways: 1) it may be applied to do work on a particle, 2) it may be transferred via bonds into a bound state, or 3) it may be emitted by a bound particle and transferred to ``nearby" particles. We capture the latter case in a second field, the emission field $\mb{F}_{emit}$, which is also coupled to particle charge.
\item Write $\mb{R} = \{\mb{r}_{1} = (x_{1}, y_{1}), \ldots, \mb{r}_{N} = (x_{N}, y_{N})\}$ for the particle positions. There are at most five possible positions for the next state of a given particle, namely, the current position plus each of the four neighboring positions. The momentum is determined by the current state and chosen position, so there are exactly five possible next states, up to Pauli exclusion. The probability distribution for the next state $\mb{R}'$ is given by the master equation
\begin{align*}
P(\mb{R}') = \sum_{\mb{S}}P(\mb{R}' \mid \mb{S})P(\mb{R} = \mb{S})
\end{align*}
which reduces to
\begin{align*}
P(\mb{R}') = P(\mb{R}' \mid \mb{R})
\end{align*}
upon fixing $\mb{R}$. So we must determine the conditional next state distribution $P(\mb{R}' \mid \mb{R})$.
\item Particles mutually interact, so the coordinate variables are not in general causally independent of one another. To make sampling tractable this requires us to factorize the conditional distribution into independent pieces. Given that particles only have access to the five nearest states at each update, and the potential only extends out to these same states, particles separated by a distance of four or more are pairwise independent. It follows that we may partition the particles into eight independent sets $I_{1}, \ldots, I_{8}$. Then Gibbs sampling over the partition allows us to sample from $P(\mb{R}' \mid \mb{R})$. When sampling the appropriate conditional distribution of an independent subset $I_{\ell}$ of particles we can sample the next state of each particle in $I_{\ell}$ independently. This reduces the problem to sampling the next state of a single particle.
\item While this method avoids treating correlated particles as independent, it does not account for all of the effects of correlation. Highly correlated particles cannot be correctly sampled separately---the marginal distribution collapses in a way that does not reflect the global distribution. This occurs most prominently when particles are in a bound state. To correct for this we run between the particle-level updates a second phase of blocked Gibbs sampling on bound states. For the purposes of efficient sampling we treat particles in a bound state as perfectly correlated during this phase, fixing their positions relative to one another. This reduces the state variables to the position of the center of mass of a bound state, as well as the orientation about its center of mass. To maintain simplicity and tractability, we elect to fix the orientation. This allows us to treat the bound state in the same manner as a particle, with five candidate positions for the next state of the center of mass. It is also enabled by the fact that rotations would have to occur in $\pi/2$ radian increments, a larger scale change than a translation by a single unit distance. This suggests that such changes are more rare and thus can be neglected.
\item We begin with the particle-level update phase. Fix a particle $i$ of charge $q$ and mass $m$ and fix a shift vector $\mb{e} \in \{(0, 0), (1, 0), (-1, 0), (0, 1), (0, -1)\}$. Write
\begin{align*}
\mb{f}_{i} = q(\mb{F}_{charge}(\mb{r}_{i}) + \mb{F}_{emit}(\mb{r}_{i})) + m\mb{F}_{mass}(\mb{r}_{i}) - \nabla_{\mb{e}} U(\mb{r}_{i}; i)
\end{align*}
for the net non-Brownian force experienced by particle $i$. Let $\mb{R}'$ be the state matching $\mb{R}$ except particle $i$ is in position $\mb{r}_{i}' = \mb{r}_{i} + \mb{e}$ with momentum
\begin{align}
\label{mom_update_eq}
\mb{p}'_{i} = (1 - \gamma \Delta t)\mb{p}_{i} + \begin{cases}
0, &\textnormal{if an emission event will occur}, \\
\mb{f}_{rand}\sqrt{\Delta t} + \mb{f}_{i} \Delta t, &\textnormal{else},
\end{cases}
\end{align}
where
\begin{align*}
\Delta t = \min\left\{{mD \over |\mb{p}_{i}|}, \sqrt{mD \over |\mb{f}_{i} - \gamma\mb{p}_{i}|}, \left(mD \over |\mb{f}_{rand}|\right)^{2/3}\right\}.
\end{align*}
This gives a time scale approximately matching the expected position transition rates. The terms of (\ref{mom_update_eq}) track the change in momentum due to the viscosity of the heat bath and the various forces. Write $\mb{p}_{i,ne}$ for the momentum vector produced in the case with no emission,
\begin{align*}
\mb{p}_{i,ne} = (1 - \gamma \Delta t)\mb{p}_{i} + \mb{f}_{rand}\sqrt{\Delta t} + \mb{f}_{i}\Delta t,
\end{align*}
and $\mb{p}_{i,nv}$ for the vector obtained when removing the damping effects of viscosity,
\begin{align}
\label{mom_nv_eq}
\mb{p}_{i,nv} = \mb{p}_{i} + \mb{f}_{rand}\sqrt{\Delta t} + \mb{f}_{i}\Delta t.
\end{align}
The change in kinetic energy of the particle system (including the emission field $\mb{F}_{emit}$) between $\mb{R}$ and $\mb{R}'$ is
\begin{align*}
\Delta K_{\mb{R}\mb{R}'} &= {1 \over 2m}(|\mb{p}_{i,ne}|^{2} - |\mb{p}_{i}|^{2}).
\end{align*}
Work done on the system is stored as kinetic or potential energy in particles until it is released into the viscous heat bath via damping effects. By default a particle travels according to its momentum. Sometimes heat is borrowed from the bath to alter the trajectory of the particle. We determine the net effect of these processes by tracking the alignment between the velocity of the particle and its change in position, an analogue for the classical expression for work as force applied along a displacement. Thus the work dissipated from the particle system in the transition from $\mb{R}$ to $\mb{R}'$ is
\begin{align*}
W_{\mb{R}\mb{R}'} &= {1 \over m}(\mb{p}_{i,ne} \cdot \mb{e}).
\end{align*} Note that this is consistent with our expression for the momentum $\mb{p}'_{i}$ because an emission event can only occur if the shift vector $\mb{e}$ is $\mb{0}$. Then the heat $\Delta Q_{\mb{R}\mb{R}'}$ released into the bath in the transition between $\mb{R}'$ and $\mb{R}$ is
\begin{align*}
\Delta Q_{\mb{R}\mb{R}'} &= W_{\mb{R}\mb{R}'} - \Delta E_{\mb{R}\mb{R}'} \\
&= W_{\mb{R}\mb{R}'} - \Delta K_{\mb{R}\mb{R}'} - \Delta U_{\mb{R}\mb{R}'} \\
&= {1 \over m}(\mb{p}_{i,ne} \cdot \mb{e}) - {1 \over 2m}(|\mb{p}_{i,ne}|^{2} - |\mb{p}_{i}|^{2}) - (U(\mb{R}') - U(\mb{R})) \\
&= {1 \over m}(\mb{p}_{i,ne} \cdot \mb{e}) - {1 \over 2m}(|\mb{p}_{i,ne}|^{2} - |\mb{p}_{i}|^{2}) - (U(\mb{r}'_{i}; i) - U(\mb{r}_{i}; i)),
\end{align*}
%%% need to account for lost kinetic energy, which would normally become internal energy; persumably immediately dissipated, so work term is wrong?
%%% different (?) problem: should it be dissipated work or work done on the system? Is work being done on the system when the particle moves *due to its momentum*? And is there any double counting of work between position and momentum?
%%% the particle is doing work against the path when it moves in the direction of its momentum? Clearly the bath or potential must do work to move it against its momentum. And is this connected to the viscosity loss?
where $\Delta E_{\mb{R}\mb{R}'}$ is the energy difference between $\mb{R}'$ and $\mb{R}$ and $\Delta U_{\mb{R}\mb{R}'}$ is the potential energy difference between $\mb{R}'$ and $\mb{R}$. Following non-equilibrium Boltzmann thermodynamics\footnote{Cite England, etc.} we take
\begin{align*}
P(\mb{R}' \mid \mb{R}) = {1 \over Z}e^{\beta\Delta Q_{\mb{R}\mb{R}'}},
\end{align*}
where $Z$ is the partition function
\begin{align*}
Z = \sum_{\mb{R}'}e^{\beta\Delta Q_{\mb{R}\mb{R}'}} = \sum_{\mb{e}}e^{\beta\Delta Q_{\mb{R}\mb{R}'}}.
\end{align*}
\item It remains to describe when emission events occur and how they determine the dynamics of the emission field $\mb{F}_{emit}$. In real physical systems energy is sometimes spontaneously emitted in the form of photons or phonons when, say, a bonded electron is excited with energies near the energy scale of the bond\footnote{Cite...}. We use an analogous criterion to determine emission events in our context. The change in kinetic energy of the particle corresponding to the emission is
\begin{align}
\label{emit_energy_eq}
E_{emit} = {1 \over 2m}(|\mb{p}_{i,nv}|^{2} - |\mb{p}_{i}|^{2}),
\end{align}
meaning all of the non-viscous work done on the particle system during the state transition is released. If $E_{emit} > 0$, particle $i$ is bound, and
\begin{align*}
{|\mb{p}_{i,ne}|^{2} \over 2m} \geqslant -\epsilon U(\mb{r}_{i}; i)
\end{align*}
for a constant $\epsilon \approx 1$, that is, the kinetic energy of particle $i$ is near the scale of its potential energy, then an emission event occurs if the particle is selected to remain in the same position (i.e. $\mb{r}'_{i} = \mb{r}_{i}$ is sampled). We require that there be a positive amount of emitted energy mainly for technical reasons, although the alternative should only occur rarely; it would imply that either the particle became bound with a greater kinetic energy than the binding energy or that it remained bound while having its potential energy reduced below its kinetic energy. The energy (\ref{emit_energy_eq}) is released as a uniform finite range energy wave via $\mb{F}_{emit}$. Every lattice site $\mb{r}$ within a distance $\delta$ of $\mb{r}_{i}$ (except for $\mb{r}_{i}$ itself) receives kinetic energy
\begin{align*}
{cE_{emit} \over |\mb{r} - \mb{r}_{i}|},
\end{align*}
where $c$ is a constant specifying the fraction of the wave's energy the particle receives and satisfying
\begin{align*}
\sum_{\delta \geqslant |\mb{r} - \mb{r}_{i}| \geqslant 1}{1 \over |\mb{r} - \mb{r}_{i}|} = {1 \over c}
\end{align*}
for energy conservation, making $c$ a simple function of $\delta$. The range parameter $\delta$ could be a small constant, reflecting fast absorption of traveling energy by the bath, or it could be set so that at the boundary the energy at a point drops below a minimum energy threshold. The energy effectively passes through or around particles, meaning it is more analogous to a wave of phonons or heat than photons. If site $\mb{r}$ does not contain a particle then the energy immediately dissipates into the bath. Otherwise, suppose particle $j$ is at point $\mb{r}$. We determine the update to the emission field by considering the case where all of the incoming energy is absorbed by particle $j$. This occurs when $\mb{r} - \mb{r}_{i}$ points in the same direction as the momentum $\mb{p}_{j}$ of particle $j$. Writing $\mb{p}_{j} = a(\mb{r} - \mb{r}_{i})$ and $b(\mb{r} - \mb{r}_{i})$ for the update vector, we require that
\begin{align*}
|a(\mb{r} - \mb{r}_{i}) + (q \Delta t)b(\mb{r} - \mb{r}_{i})|^{2} - |a(\mb{r} - \mb{r}_{i})|^{2} = {2m_{j}cE_{emit} \over |\mb{r} - \mb{r}_{i}|}.
\end{align*}
This can be readily solved, yielding
\begin{align*}
b = {1 \over q_{j} \Delta t}\left(-a + \sqrt{a^{2} + {2m_{j}cE_{emit} \over |\mb{r} - \mb{r}_{i}|^{3}}}\right).
\end{align*}
Substituting $a = |\mb{p}_{j}| / |\mb{r} - \mb{r}_{i}|$ and simplifying, we update the emission field by setting
\begin{align*}
\mb{F}'_{emit}(\mb{r}) = \mb{F}_{emit}(\mb{r}) + {\mb{r} - \mb{r}_{i} \over q_{j} \Delta t|\mb{r} - \mb{r}_{i}|}\left(\sqrt{|\mb{p}_{j}|^{2} + {2m_{j}cE_{emit} \over |\mb{r} - \mb{r}_{i}|}} - |\mb{p}_{j}|\right).
\end{align*}
\item The second phase of Gibbs sampling over bound states works the same way as Gibbs sampling over particles with minor changes. The mass of a bound state $A$ is
\begin{align*}
m_{A} = \sum_{a \in A}m_{a},
\end{align*}
where $m_{a}$ is the mass of particle $a$, and the position of $A$ is its center of mass,
\begin{align*}
\mb{r}_{A} = {1 \over m_{A}}\sum_{a \in A}m_{a}\mb{r}_{a}.
\end{align*}
The potential energy and momentum of $A$ are simply determined additively,
\begin{align*}
U(A) &= \sum_{a \in A}U(\mb{r}_{a}; a), \\
\mb{p}_{A} &= \sum_{a \in A}\mb{p}_{a},
\end{align*}
and substitute for the particle equivalents as expected. There are no emission events for bound states.
\item Since the probability distribution $P(\mb{R}' \mid \mb{R})$ is composed of only five terms for both phases, we can tractably sample from it by computing the full distribution exactly.
\end{itemize}

%{1 \over \mu q_{j}}\left((\mb{n} \cdot \mb{p}_{j}) + \sqrt{(\mb{n} \cdot \mb{p}_{j})^{2} + {cE_{emit} \over |\mb{r} - \mb{r}_{i}|}}\right)


% Examining (\ref{emit_energy_eq}) and (\ref{mom_nv_eq}), we update the emission field by setting
% \begin{align*}
% \mb{F}'_{emit}(\mb{r}) = \mb{F}_{emit}(\mb{r}) + \mb{w},
% \end{align*}
% where the received field vector $\mb{w}$ satisfies the system
% \begin{align*}
% |\mb{p}_{j} + \mu q\mb{w}|^{2} - |\mb{p}_{j}|^{2} &= {2m_{j}cE_{emit} \over |\mb{r} - \mb{r}_{i}|} \\
% {\mb{w} \over |\mb{w}|} &= \mb{n}
% \end{align*}
% for $m_{j}$ the mass of particle $j$ and $\mb{n}$ the unique unit vector in the direction of $\mb{r} - \mb{r}_{i}$. Clearly this ensures that particle $j$ receives exactly $cE_{emit}/|\mb{r} - \mb{r}_{i}|$ units of kinetic energy. Furthermore, there exists a solution $\mb{w}$ because the right-hand side of the first equation is positive. We can combine and simplify the system of equations to obtain
% \begin{align*}
% |\mb{p}_{j} + s\mu q\mb{n}| = \sqrt{|\mb{p}_{j}|^{2} + {2m_{j}cE_{emit} \over |\mb{r} - \mb{r}_{i}|}}
% \end{align*}
% and it suffices to solve for $s > 0$. This can be solved piecewise. Writing $A$ for the right-hand side of the above equation, $\mb{p}_{j} = (p_{j,1}, p_{j,2})$, $\mu q\mb{n} = (n_{1}, n_{2})$, $\mb{p}_{j} + s\mu q\mb{n} = (v_{1}, v_{2})$, $\sigma_{1} = \sign(v_{1})$, and $\sigma_{2} = \sign(v_{2})$, $s_{1} = -p_{j,1}/n_{1}$, and $s_{2} = -p_{j,2}/n_{2}$, we have
% \begin{align*}
%  \eta s = A - \theta,
% \end{align*}
% where $\eta = \sigma_{1}n_{1} + \sigma_{2}n_{2}$ and $\theta = \sigma_{1}p_{j,1} + \sigma_{2}p_{j,2}$. Then
% \begin{align*}
% ...
% \end{align*}

Draft model idea 5:
\begin{itemize}
\item Coordinate space is a cylindrical lattice $[n] \times [n]$ for $n \in \mathbb{N}$, with a periodic boundary on the first dimension; we write the points as vectors $\mb{r} = (x, y)$.
\item The space is populated by $N$ particles of $k$ different types of counts $n_{1}, \ldots, n_{k}$, $\sum n_{i} = N$, with positional Pauli exclusion amongst all particles. Particles of a common type are indistinguishable.
\item Particle types $1, \ldots, k$ have masses $m_{1}, \ldots, m_{k}$ and charges $-1, 1, 2, \ldots, k-1$, denoted by $q_{1}, \ldots, q_{k}$. Each particle $i$ has a position $\mb{r}_{i}$ and momentum $\mb{p}_{i}$. The kinetic energy $K(i)$ of particle $i$ given by $|\mb{p}_{i}|^{2}/2m_{i}$, where $|\mb{v}| = |x| + |y|$ is the 1-norm of $\mb{v} = (x, y)$.
\item Particles interact via a potential and a bonding mechanism. Two distinct particles $i$ and $j$ have joint potential energy
\begin{align*}
U_{ij} = \begin{cases}
q_{i}q_{j} + {1 \over 2}(k_{1}\sigma_{1,ij}^{2} + k_{2}\sigma_{2,ij}^{2}), &b_{ij} = 1, \\
\eta q_{i}q_{j}, &b_{ij} = 0 \textrm{ and } |\mb{r}_{i} - \mb{r}_{j}| = 1, \\
0, &\textrm{else},
\end{cases}
\end{align*}
where $0 < \eta \ll 1$ is a constant tuning the strength of the potential, $b_{ij}$ is the bond indicator variable, $\sigma_{1,ij}$ and $\sigma_{2,ij}$ are the longitudinal and transverse bond stresses (parallel and perpendicular, respectively, to the bond), and $k_{1}$ and $k_{2}$ are constants tuning how energy is stored as stress in a bond. All particles weakly interact at close range, with the polarity of the interaction determined by their charges. Stronger interactions are mediated by binary bonding, which is active for particles $i$ and $j$ when $b_{ij} = 1$. We essentially treat each bond as a small harmonic oscillator in a potential well whose height is proportional to the charges of the bonded particles. We require that particles be adjacent in order to bond, that is, $b_{ij} = 0$ if $|\mb{r}_{i} - \mb{r}_{j}| > 1$. Bonding is meant to be the overwhelmingly dominant mechanism causing particles to become bound, so the constant $\eta$ should be set in a way that ensures nonbonding interactions generically only lead to short-lived correlations between particles. The potential energy $U(i)$ of particle $i$ is determined additively:
\begin{align*}
U(i) = \sum_{j}U_{ij}.
\end{align*}
\item We say that a particle $i$ is {\em bound} to each of its neighbors with which it shares a bond of negative energy. When this relation is extended to an equivalence relation, we call the equivalence classes {\em bound states}. Equivalently, these are the connected components in the graph generated by connecting particles with an edge if they share a negative energy bond.
\item The particles are immersed in a viscous heat bath with temperature $T = 1/\beta$ and collision coefficient $\gamma$. Viscosity introduces both damping and fluctuations, with the latter appearing as a Brownian force
\begin{align*}
\mb{f}_{rand} \sim \sqrt{2\gamma m \over \beta}\boldsymbol{\eta},
\end{align*}
where $\boldsymbol{\eta}$ is a 2-dimensional standard normal distribution and $m$ is the mass of a particle subject to the force.
\item Adjacent particles can shield a particle from a portion of the viscosity of the bath. We take $\gamma$ to be the maximum viscosity experienced by a particle and $\gamma_{min}$ to be the minimum. For a lattice site $\mb{r}$, let $\epsilon(\mb{r})$ be the fraction of adjacent lattice sites that are unoccupied by particles. Then
\begin{align*}
\gamma(\mb{r}) = \gamma_{min} + (\gamma - \gamma_{min}){3\epsilon(\mb{r}) \over 2\epsilon(\mb{r}) + 1}
\end{align*}
gives the viscosity experienced by particle at $\mb{r}$. The residual shielding of an additional adjacent particle increases as the number of adjacent particles increases. This shielding mechanism enables state-dependent differences in energy dissipation rates and reflects the general phenomena that dissipation involves the dispersal of energy. Particles will often move between adjacent lattice sites. During any such transition we set the viscosity experienced by the particle to the arithmetic mean of the viscosities at each site.
\item Changes in particle state can transform potential energy into kinetic energy and vice versa. We now define the forces delivering that transformation. For a continuous scalar field $P$ of potential energies the force due to $P$ is given by the negative gradient of $P$, $-\nabla P$. The state space considered here is discrete, so we require a different route to defining the force produced by potential energy differences. We first note that we require an expression for the interaction force on a particle $i$ in order to determine the change in momentum under a change in the potential energy of $i$. Suppose then that the system state changes and this causes a change in the potential energy of $i$. By additivity we may assume that the change is caused by the interaction of $i$ with exactly one other particle $j$. It is natural to define a directional gradient $\nabla U(i)$ of the potential energy of particle $i$ as the vector satisfying the equation
\begin{align*}
{U'_{ij} - U_{ij} \over 2} = \nabla U(i) \cdot \mb{e}
\end{align*}
%%% consider each neighbor pair, i.e. each interaction separately
for a unit vector $\mb{e}$ specifying the direction of $\nabla U$, where $U'_{ij}$ is the interaction energy of $i$ and $j$ after the state change. Essentially each particle receives half the potential energy contained in their interaction. This gives the correct expression for the work done by the interaction force and mirrors one standard definition of the gradient via the total derivative.
It also implies that
\begin{align*}
\nabla U(i) = {U'_{ij} - U_{ij} \over 2}\mb{e}.
\end{align*}
To specify $\mb{e}$ there are two cases to consider: (1) the state change is due to a change in particle positions, and (2) it is due to a change in bond states. For the former we set
\begin{align*}
\mb{e} = \begin{cases}
\mb{r}_{i}' - \mb{r}_{i}, &i \textrm{ moved}, \\
\mb{r}_{j} - \mb{r}_{j}', &\textrm{else,}
\end{cases}
\end{align*}
where $\mb{r}_{i}'$ and $\mb{r}_{j}'$ are the new positions of particle $i$ and $j$, respectively. We cannot employ the same strategy in case (2) because there is no change in position, but a similar approach can be employed. Let $\mb{e}_{ij} = \mb{r}_{i} - \mb{r}_{j}$ and $\mb{e}^{\perp}_{ij}$ be the unit vector anticlockwise perpendicular to $\mb{e}_{ij}$. Then we take
\begin{align*}
\mb{e} = {|\sigma_{1,ij}|\mb{e}_{ij} + \sigma_{2,ij}\mb{e}^{\perp}_{ij} \over |\sigma_{1,ij}| + |\sigma_{2,ij}|}.
\end{align*}
This captures the behavior of an asymmetric two-dimensional harmonic oscillator.
\item The particles are subject to up to two space-dependent external drives represented by fields $\mb{F}_{charge}$ and $\mb{F}_{mass}$ coupled to particle charge and mass, respectively. That is, the force generated by each field on a particle of charge $q$ and mass $m$ is $q\mb{F}_{charge}$ and $m\mb{F}_{mass}$, respectively. In general these drives could take any form, whether stochastic or deterministic, and interesting behavior might emerge across a variety of drive choices. We choose to focus on physically-inspired drives that evoke the spectrum of solar radiation near Earth as well as its gravitational field, capturing the dominant energy sources on Earth. The charge-coupled field is modeled after electromagnetism. We write
\begin{align*}
\mb{F}_{charge} = \mb{E} + {\mb{p} \over \kappa m} \times \mb{B},
\end{align*}
where $\mb{p}$ is the momentum of a particle influenced by $\mb{F}_{charge}$ and $\kappa$ is a constant determining the relative strength of $\mb{B}$ and $\mb{E}$. Assuming $\mb{E}$ and $\mb{B}$ are generated by blackbody radiation traveling vertically in from the top of the lattice, in accordance with Planck's law we take
\begin{align*}
\mb{E}((x, n)) &= (\theta, 0), \\
\theta &\sim {x^{3} \over e^{\alpha x} - 1}, \textrm{ and} \\
\mb{E}((x, y)) &= \psi((x, y))\mb{E}((x, n)),
\end{align*}
where $\alpha$ is a constant tuning the strength of the drive and $\psi$ specifies how the drive varies in space. We define $\mb{B}$ as the field orthogonal to both coordinate dimensions of the lattice and of equal magnitude to $\mb{E}$ such that $\mb{E} \times \mb{B}$ points directly downwards in the vertical axis, i.e. in the direction of the radiation. To emulate solar radiation near Earth, the constant $\alpha$ should be set such that the maximum of the distribution of $\theta$ is around an order of magnitude less than the strongest binding energies attained by particles in the system. Radiation cannot in general pass through particles unhindered. Therefore $\psi$ should vary across the lattice in a manner that reflects this screening effect. We could treat particles as perfect barriers or absorbers and set
\begin{align*}
\psi((x, y)) = \begin{cases}
0, &\textrm{there exists a particle in } \{x\} \times (y, n], \\
1, &\textrm{else.}
\end{cases}
\end{align*}
We could also instead consider particles as imperfect absorbers and assume an exponential falloff in intensity, setting
\begin{align*}
\psi((x, y)) = \mu_{1}e^{-\mu_{2}\ell} \textrm{ if there exist } \ell > 0 \textrm{ particles in } \{x\} \times (y, n].
\end{align*}
The constants $\mu_{1}$ and $\mu_{2}$ should be set such that an appropriate normalization holds:
\begin{align*}
\sum_{\ell = 0}^{\infty}\mu_{1}e^{-\mu_{2}\ell} = 1,
\end{align*}
giving $\mu_{1} = 1 - e^{\mu_{2}}$. The mass-coupled field $\mb{F}_{mass}$ is a uniform field modeled after gravitation. Fixing a constant $g$, we set
\begin{align*}
\mb{F}_{mass}(\mb{r}) = g
\end{align*}
everywhere. The constant $g$ should be set such that in general $q\mb{F}_{charge}$ dwarfs $m\mb{F}_{mass}$ unless $m$ is much larger than $q$, although one could also consider alternative schemes.
\item We use $\tau$-leaping, a variant of the kinetic Monte Carlo or Gillespie algorithm, to simulate the system. Write $\mb{X}(t) = \{\mb{R}(t), \mb{P}(t), \mb{B}(t), \boldsymbol{\Sigma}(t)\}$ for the state of the system at time $t$. The dynamics of the system are determined by elementary events $e: \mb{X} \rightarrow \mb{X}'$ that may occur at any given time. These events take place at state-dependent rates $r_{e}(\mb{X})$. There are five types of events:
\begin{itemize}
\item[--] Particle hop $(i, \mb{e})$: particle $i$ moves from $\mb{r}_{i}$ to $\mb{r}_{i} + \mb{e}$ for a shift vector $\mb{e} \in \{(1, 0), (-1, 0), (0, 1), (0, -1)\}$.
\item[--] Elastic collision $(i, \mb{e})$: particle $i$ interacts with and transfers energy to particles adjacent to $\mb{r}_{i} + \mb{e}$ in the manner of an elastic collision, remaining in the same position.
\item[--] Bond flip $(i, j)$: adjacent particles $i$ and $j$ change bond state from $b_{ij}$ to $1 - b_{ij}$.
\item[--] Stress relaxation: bond stress diffuses between incident bonds.
\item[--] Bound state hop $(A, \mb{e})$: bound state $A$ moves from $\mb{r}_{A}$ to $\mb{r}_{A} + \mb{e}$ for a shift vector $\mb{e}$, where $\mb{r}_{A}$ is the center of mass of $A$.
\end{itemize}
We treat bond stress relaxation as occurring at much faster time scales than the other events. This enables us to perform the event deterministically between the other events, simplifying the dynamics and reducing sampling costs. Every other event has an associated rate: $r^{H}_{i,\mb{e}}$, $r^{C}_{i,\mb{e}}$, $r^{F}_{i,j}$, and $r^{H}_{A,\mb{e}}$, respectively.
\item All four (stochastic) event types can mutually effect each other. As a result we run four separate phases of the sampling process, one for each event type. Furthermore, particles mutually interact, so their movement events are not in general causally independent of one another. This requires us to factorize the set of particles into independent subsets. Given that particles only have access to the four nearest lattice sites at each update, and both interaction mechanisms only extend out to these same positions, particles separated by a distance of four or more are independent during a single hop. It follows that we may partition the particles into eight independent sets $I_{1}, \ldots, I_{8}$. The same considerations hold for bound states, with a need to partition them into a similar number of independent sets. Sampling occurs in series across these sets and in parallel within them.
\item Letting
\begin{align*}
\Lambda^{H}_{i} = \sum_{\mb{e}}r^{H}_{i,\mb{e}}(\mb{X}(t')), \quad \Lambda^{C}_{i} = \sum_{\mb{e}}r^{C}_{i,\mb{e}}(\mb{X}(t')), \quad \Lambda^{F}_{i,j} = r^{F}_{i,j}(\mb{X}(t')), \quad \Lambda^{H}_{A} = \sum_{\mb{e}}r^{H}_{A,\mb{e}}(\mb{X}(t'))
\end{align*}
for $t' < t$ the time at the last sampling step, define
\begin{align*}
\tau = {\theta \over \max\{\Lambda^{H}_{i}, \Lambda^{C}_{i}, \Lambda^{F}_{i,j}, \Lambda^{H}_{A}\}}
\end{align*}
for a constant $\theta \leqslant 1$. This determines the effective time interval in which we are sampling events; it is chosen to approximately ensure that each object---particle, bond, or bound state---experiences at most one of each event type in the time interval. For each object-event type pair $(a, \mc{E})$, the probability $P(a, \mc{E})$ that it occurs in a time interval of length $\tau$ is given by the Poisson CDF
\begin{align*}
P(a, \mc{E}) = 1 - e^{-\Lambda^{\mc{E}}_{a}\tau},
\end{align*}
where $\Lambda^{\mc{E}}_{a}$ is now computed using the current time $t$. Then we sample from the corresponding Bernoulli distribution to determine which objects undergo a given event type in the interval $[t, t + \tau)$. For objects $a$ that undergo an event type $\mc{E}$, we sample from the categorical distribution
\begin{align*}
P(e \in \mc{E}) = {r_{a,e} \over \Lambda^{\mc{E}}_{a}}
\end{align*}
to determine which event in $\mc{E}$ occurs for $a$. Doing this across all objects and event types as described above produces the state $\mb{X}(t + \tau)$ of the system at time $t + \tau$.
\item We now begin to describe the various events, starting with stress relaxation, which is updated before the bond flip phase. The next bond stresses only depend on the current state, simplifying the dynamics. Bonds are stressed when either there is an imbalance in the particle momenta across the bond or when adjacent bonds due to their own stress differentially pull on the particles within the bond. We begin by describing the first mechanism, the source of bond stress. Suppose particles $i$ and $j$ are bonded. Writing $\mb{e}_{ij} = \mb{r}_{i} - \mb{r}_{j}$ and $\mb{e}^{\perp}_{ij}$ for the unit vector anticlockwise perpendicular to $\mb{e}_{ij}$, following Hooke's law we set
\begin{align}
\label{stress_source_eq}
s_{1,ij} = {\mb{e}_{ij} \over \sqrt{k_{1}}} \cdot {\mb{p}_{i} - \mb{p}_{j} \over \sqrt{m}} \quad \textrm{and} \quad s_{2,ij} = {\mb{e}^{\perp}_{ij} \over \sqrt{k_{2}}} \cdot {\mb{p}_{i} - \mb{p}_{j} \over \sqrt{m}}
\end{align}
These expressions are chosen to transform kinetic energy from the particles into potential energy in the bond. The second mechanism essentially diffuses stresses between neighboring bonds. Longitudinal stress along a bond is transferred to longitudinal stress in parallel neighbor bonds and transverse stress in perpendicular neighbor bonds, while transverse stress across a bond is transferred to transverse stress in parallel neighbor bonds and longitudinal stress in perpendicular neighbor bonds. Indexing the bonds around a particle $i$ in each of the compass directions with $(i, N), (i, E), (i, S), (i, W)$, write
\begin{align*}
(\nabla \cdot \sigma)_{x}(i) &= \sigma_{2}(i, N) - \sigma_{2}(i, S) - \sigma_{1}(i, E) + \sigma_{1}(i, W), \textrm{ and} \\
(\nabla \cdot \sigma)_{y}(i) &= \sigma_{1}(i, N) - \sigma_{1}(i, S) + \sigma_{2}(i, E) - \sigma_{2}(i, W),
\end{align*}
the bond stress divergence in each dimension at $i$, where $\sigma_{\cdot}(i, \cdot) = 0$ for each missing bond. Then these obey the force balance equations
\begin{align*}
(\nabla \cdot \sigma)_{x}(i) &= (\nabla \cdot s)_{x}(i), \textrm{ and} \\
(\nabla \cdot \sigma)_{y}(i) &= (\nabla \cdot s)_{y}(i).
\end{align*}
This is essentially a discretized version of Laplace's equation over bond stress divergence at the particle sites. With this we can define the bond dynamics phase, which occurs prior to particle and bound state sampling. We first calculate (\ref{stress_source_eq}) for every bond and set
\begin{align*}
\boldsymbol{\sigma}_{ij} = (1 - \gamma)\boldsymbol{\sigma}_{ij} + \mb{s}_{ij},
\end{align*}
indicating that bonds accumulate stress from kinetic sources and lose stress due to viscous damping from the heat bath. Force balance is then enforced using an iterative smoothing method. We calculate the bond stress divergences at each particle, average over the nearest neighbor divergences, and then assign each bond the average of value of these smoothed divergences. For computational efficiency we perform a small number $i \approx 3$ of smoothing iterations. This choice also reflects the relative speeds of particles and sound on Earth.
\item For an elementary event $e: \mb{X} \rightarrow \mb{X}'$, let $e^{*}: \hat{\mb{X}}' \rightarrow \hat{\mb{X}}$ be the event ``inverse" to $e$, that is, the time-reversed transition obtained by negating all momenta. Then nonequilibrium thermodynamics requires that
\begin{align*}
\log{r_{e}(\mb{X}) \over r_{e^{*}}(\mb{X}')} = -\beta[E(\mb{X}') - E(\mb{X})] + \beta W_{e}(\mb{X}),
\end{align*}
where $E(\mb{X})$ is the energy of $\mb{X}$ and $W_{e}(\mb{X})$ is the work done on the system during $e$. This is a consequence of Crooks' fluctuation theorem and can be cast as a local detailed balance condition. Taking
\begin{align*}
r_{e}(\mb{X}) = \nu_{e}e^{\beta[W_{e}(\mb{X}) - \Delta E_{e}(\mb{X})]/2},
\end{align*}
where $\Delta E_{e}(\mb{X}) = E(\mb{X}') - E(\mb{X})$ and $\nu_{e} = \nu_{e^{*}}$ is an event-symmetric constant, gives rates satisfying local detailed balance. The energy of the system is given by the sum of the kinetic and potential energies,
\begin{align*}
E(\mb{X}) = K(\mb{X}) + U(\mb{X}),
\end{align*}
which are
\begin{align*}
K(\mb{X}) = \sum_{i}{|\mb{p}_{i}|^{2} \over 2m_{i}}
\end{align*}
and
\begin{align*}
U(\mb{X}) = \sum_{\substack{i,j \\ |\mb{r}_{i} - \mb{r}_{j}| = 1}}U_{ij} = \sum_{\substack{i,j \\ |\mb{r}_{i} - \mb{r}_{j}| = 1}}[\eta + b_{ij}(1 - \eta)]q_{i}q_{j} + {b_{ij} \over 2}(k_{1}\sigma_{1,ij}^{2} + k_{2}\sigma_{2,ij}^{2}).
\end{align*}
The work done on the system during an event $e$ is
\begin{align*}
W_{e}(\mb{X}) = \int_{C_{e}}\mb{f}_{ext} \cdot d\mb{s} \approx \mb{f}_{ext} \cdot (\mb{R}' - \mb{R}),
\end{align*}
where the approximation is obtained by considering the external force $\mb{f}_{ext}$ to be constant during the event. It remains only to determine $\nu_{e}$. The event space does not vary by momenta, indicating that we only consider a single value for the new momentum of each particle. We need to ensure that we properly account for the distribution of possible next momenta. The particles take on underdamped Langevin dynamics through the coupling with the viscous heat bath. It is known that the transition distribution over time $\tau$ for the momenta of a particle of mass $m$ and momentum $\mb{p}$ obeying Langevin dynamics is approximately given by a normal distribution of mean
\begin{align*}
\mb{p} + \tau(\mb{f} - \gamma\mb{p})
\end{align*}
and variance
\begin{align*}
{2m\gamma\tau \over \beta},
\end{align*}
where $\mb{f}$ is the non-drag force on the particle. Writing $T_{\tau}(\mb{p} \rightarrow \mb{p}')$ for this transition distribution, the nonsymmetric part is given by
\begin{align*}
{T_{\tau}(\mb{p} \rightarrow \mb{p}') \over T_{\tau}(-\mb{p}' \rightarrow -\mb{p})} = e^{\beta\Delta K_{e}(\mb{X})}.
\end{align*}
Thus by our choice of $r_{e}$ we need only consider the symmetric part,
\begin{align*}
[T_{\tau}(\mb{p} \rightarrow \mb{p}') T_{\tau}(-\mb{p}' \rightarrow -\mb{p})]^{1/2}.
\end{align*}
We always select the mean for the new momentum of a particle, meaning that this expression reduces to the leading coefficient of $T_{\tau}(\mb{p} \rightarrow \mb{p}')$:
\begin{align*}
\sqrt{\beta \over 4\pi m\gamma\tau}.
\end{align*}
This gives the first factor of $\nu_{e}$. Now consider the limiting case of a single isolated particle traveling with nonzero momentum in a bath of zero viscosity, zero temperature, and no external drive. In this scenario the next position of the particle should be totally determined by its momentum. Since the Boltzmann factor is always 1 in this case, the mobility induced by the momentum is missing from the expression for the rate $r_{e}$. In particular we are missing the effects of momentum on the change of position of a particle. We can account for this by factoring in the probability that a particle jumps at least one lattice site. Fix a particle of mass $m$ and momentum $\mb{p}$ at time $0$. Writing
\begin{align*}
\mb{r}_{\tau, \mb{e}} = \int_{0}^{\tau}\mb{e} \cdot {\mb{p}(t) \over m}dt
\end{align*}
for the displacement towards the unit vector $\mb{e}$ during $[0, \tau]$, define
\begin{align*}
\lambda_{\tau}(\mb{p}) = {1 \over \tau}P(\mb{r}_{\tau, \mb{e}} \geqslant D),
\end{align*}
where $D$ is a constant specifying the distance between adjacent lattice sites. This is the crossing rate for the particle in the direction of $\mb{e}$. Solving the Langevin equation gives
\begin{align*}
P(\mb{r}_{\tau, \mb{e}} \geqslant D) = 1 - \Phi\left({D - \mu \over \sigma}\right),
\end{align*}
where $\Phi$ is the standard normal CDF,
\begin{align*}
\mu = {f_{\mb{e}} \over m\gamma}\tau + \left({p_{\mb{e}} \over {m}} - {f_{\mb{e}} \over m\gamma}\right){1 - e^{-\gamma\tau} \over \gamma},
\end{align*}
and
\begin{align*}
\sigma^{2} = {1 \over \beta m\gamma^{2}}\left(2\gamma\tau - 4(1 - e^{-\gamma\tau}) + (1 - e^{-2\gamma\tau})\right).
\end{align*}
We have written $v_{\mb{e}}$ for $\mb{e} \cdot \mb{v}$. Note also that we have used $\gamma$ for the viscosity for brevity---it should actually be replaced by the appropriate state-dependent, transition-averaged value for $\gamma$ as defined previously. Then $\nu_{e}$ is the product of the momentum correction and the crossing rate. This concludes the specification of the rates for events involving a change of total energy and position. It remains to define the events themselves, as well as to specify the rates for events that do not involve a change in total energy or position---this includes inelastic collisions and bond flips.
%%% approx?
\item Fix a particle $i$ of charge $q$ and mass $m$ and a hop event $e$ involving particle $i$. The external force experienced by particle $i$ is
\begin{align*}
\mb{f}_{ext,i} = q\mb{F}_{charge}(\mb{r}_{i}) + m\mb{F}_{mass}(\mb{r}_{i}).
\end{align*}
while the net force is given by
\begin{align*}
\mb{f}_{i} = \mb{f}_{ext,i} + \mb{f}_{rand} - \nabla_{e}U(i) - \gamma\mb{p}_{i},
\end{align*}
where $\nabla_{e}U(i)$ is the interaction force on $i$ under event $e$. Then particle $i$ attains momentum
\begin{align*}
\mb{p}'_{i} = \mb{p}_{i} + \tau\mb{f}_{i}
\end{align*}
under event $e$, while every other particle $j$ attains momentum
\begin{align*}
\mb{p}_{j}' = \mb{p}_{j} - \tau\nabla_{e}U(j).
\end{align*}
\item Bound state hops work the same way as particle hops with minor differences. The mass of a bound state $A$ is
\begin{align*}
m_{A} = \sum_{a \in A}m_{a},
\end{align*}
where $m_{a}$ is the mass of particle $a$, and the position of $A$ is its center of mass,
\begin{align*}
\mb{r}_{A} = {1 \over m_{A}}\sum_{a \in A}m_{a}\mb{r}_{a}.
\end{align*}
The potential energy and momentum of $A$ are simply determined additively,
\begin{align*}
U(A) &= \sum_{a \in A}U(\mb{r}_{a}; a), \\
\mb{p}_{A} &= \sum_{a \in A}\mb{p}_{a},
\end{align*}
and substitute for the particle equivalents as expected.
\item We have already given sufficient information to specify every event type other than elastic collisions and bond flips. We begin with the former. Fix a particle $i$ and step vector $\mb{e} \in \{(1, 0), (-1, 0), (0, 1), (0, -1)\}$. Writing $N(\mb{r})$ for the particles neighbouring $\mb{r}$, let
\begin{align*}
N = N(\mb{r}_{i}) \cup N(\mb{r}_{i} + \mb{e}) - \{i\}
\end{align*}
and
\begin{align*}
\mb{g} = \sum_{j \in N}\eta q_{i}q_{j}(\mb{r}_{j} - \mb{r}_{i})
\end{align*}
be a gradient vector. Then particle $i$ attains momentum
\begin{align*}
\mb{p}'_{i} = \mb{p}_{i} - J\mb{n},
\end{align*}
where
\begin{align*}
\mb{n} = \begin{cases}
\mb{e}, &\textrm{if } \mb{r}_{i} + \mb{e} \textrm{ is occupied}, \\
\mb{g}/(|\mb{g}| + \omega), &\textrm{else.}
\end{cases}
\end{align*}
for a small constant $\omega$ and
\begin{align*}
J = {2m_{i}M \over m_{i} + M}u
\end{align*}
for $u = \left(\mb{p}_{i}/m_{i} - V\right) \cdot \mb{n}$. The terms $M$ and $V$ are the effective mass and velocity of the particles $i$ is reacting against, namely those contained in the set
\begin{align*}
S = \begin{cases}
j, &\textrm{if } \mb{r}_{j} = \mb{r}_{i} + \mb{e}, \\
N(\mb{r}_{i} + \mb{e}), &\textrm{else.}
\end{cases}
\end{align*}
For $j \in  S$, write
\begin{align*}
w'_{j} = \left|q_{i}q_{j}(\mb{e}_{ij} \cdot \mb{n})\right| \quad \textrm{and} \quad w_{j} = {w'_{j} \over \sum_{j}w'_{j}},
\end{align*}
where $\mb{e}_{ij}$ is the unit vector form $i$ toward $j$. Then
\begin{align*}
M = \sum_{j \in S}w_{j}m_{j} \quad \textrm{and} \quad V = {1 \over M}\sum_{j \in S}w_{j}\mb{p}_{j},
\end{align*}
and each particle $j \in S$ attains momentum
\begin{align*}
\mb{p}'_{j} = \mb{p}_{j} + {w_{j}m_{j} \over M}J\mb{n}.
\end{align*}
Let $\Delta U_{\mb{e}}$ be the energy difference produced by moving particle $i$ from $\mb{r}_{i}$ to $\mb{r}_{i} + \mb{e}$ and
\begin{align*}
K_{\mb{e}} = {1 \over 2m_{i}}\max\{\mb{p}_{i} \cdot \mb{e}, 0\}^{2}
\end{align*}
be the kinetic energy of $i$ in the direction of $\mb{e}$. Then the rate $r^{C}_{i,\mb{e}}$ of the event is given by
\begin{align*}
r^{C}_{i,\mb{e}}(\mb{X}) &= \nu_{C}\mb{1}[K_{n} > 0 \textrm{ and } i \textrm{ is not bound} \textrm{ and } (r_{i} + \mb{e} \textrm{ is occupied or } K_{n} < \Delta U_{\mb{e}})] \\
&\qquad\times \sigma(\rho(\max\{u, 0\} - u_{0})),
\end{align*}
where $\sigma$ is the sigmoid function, $\rho$ is a steepness constant, $u_{0}$ is threshold velocity scale, and $\nu_{C}$ should be chosen to set the base rate to the same scale as other event rates.
\item
\item Since all probability distributions are composed of five or fewer terms, we can tractably sample from them by computing the full distribution exactly.
\end{itemize}


\subsection{Back to the drawing board}

Starting system framework: units\footnote{Basic element, particle, etc. usually corresponding to the minimal independent sets of state variables} fluctuating between states interacting with other units. Units are immersed in a (possibly viscous) heat bath. Unit interactions are characterized by an interaction or potential energy determined by their states. State transformations depend on unit energies (which are the addition of potential energy and other energies) and external potentials. They may transform energy between forms, transport energy across units, or exchange energy with the bath.

The goal: produce a system that transports energy from external potentials into the bath and through that transportation self-organizes to make the process more efficient.

Which units interact?
\begin{itemize}
\item TSMN: small geometrically fixed set for each unit
\item Pure particle model: Units have lattice position as a state variable, neighboring units interact
\end{itemize}

What energy transformations are possible? All energy gradients should eventually be thermalized.
\begin{itemize}
\item To what extent is the ``dissipation via dispersion" phenomenon an enabler of open-endedness? This seems to combine with the dispersed nature of excitations as a dominant cause of open-endedness
\end{itemize}
High-level idea: units assemble in shapes/graphs with geometric or configurational energy barriers that can only be navigated or overcome by complementarily shaped assemblies.
\begin{itemize}
\item Could setup for variable number of neighbors to allow for less uniform geometries
\end{itemize}
It seems that many of the core features of our particle approach essentially follow from this idea, in tandem with \textit{naturalness}\footnote{In particular an energy-based framework and no highly ad hoc or artificial mechanisms} and the desire to keep system complexity outside of the individual units (i.e. instead in their interactions). What remains to be refined or fundamentally changed is the energy dynamics. Could it make sense to remove momentum or simplify it into something like temperature? Where can memory happen with such a change?
\begin{itemize}
\item There seems to be a sort of incompatibility or tension between momentum as a full continuous degree of freedom and a discrete, short range potential, especially a simple binary interaction. Something else that is simpler is needed.
\end{itemize}

The fundamental hypothesis is a scale hypothesis: a thermodynamic system with few dynamical ingredients can produce complex open-ended behavior at large scales.

New idea:
\begin{itemize}
\item position and temperature / kinetic energy are the degrees of freedom for units
\item universal weak short-range attractive interaction between units
\item bonding mechanism with discrete heat diffusion across bonds
\item two events: bound state hops and bond flips
\end{itemize}
This foregoes polar or typic energy barriers, leaving only configurational or geometric energy barriers. In order to bond, two bound states need to have sufficiently many units neighboring in order to overcome their thermal and configurational kinetic energy; e.g. 1 neighbor is generally nontrapping, 2 neighbors is trapping for short times, and 3 neighbors is trapping for long times.

Draft model idea 6:
\begin{itemize}
\item Coordinate space is a cylindrical lattice $[n] \times [n]$ for $n \in \mathbb{N}$, with a periodic boundary on the first dimension; we write the points as vectors $r = (x, y)$.\footnote{This is the default choice for concreteness and to mirror Earth's surface, but it could be any lattice.}
\item The space is populated by $N$ indistinguishable particles, with positional Pauli exclusion amongst all particles.
\item All particles have mass $m$. Each particle $i$ has a position $\mb{r}_{i}$ and internal energy $E_{i}$.
\item Particles interact via a potential and a bonding mechanism. Two distinct particles $i$ and $j$ have joint potential energy
\begin{align*}
U_{ij} = \begin{cases}
-1 + k\sigma_{ij}^{2}, &b_{ij} = 1, \\
-\eta, &b_{ij} = 0 \textnormal{ and } |\mb{r}_{i} - \mb{r}_{j}| = 1, \\
0, &\textnormal{else},
\end{cases}
\end{align*}
where $0 < \eta \ll 1$ is a constant tuning the strength of the potential, $b_{ij}$ is the bond indicator variable, $\sigma_{ij}$ is the bond stress, and $k$ is a constant tuning how energy is stored as stress in a bond. All particles interact weakly and attractively at close range. Stronger attractive interactions are mediated by binary bonding, which is active for particles $i$ and $j$ when $b_{ij} = 1$. We essentially treat each bond as a small harmonic oscillator in a potential well with constant height. We require that particles be adjacent in order to bond, that is, $b_{ij} = 0$ if $|\mb{r}_{i} - \mb{r}_{j}| > 1$. Bonding is meant to be the overwhelmingly dominant mechanism causing particles to become bound, so the constant $\eta$ should be set in a way that ensures non-bonding interactions generically only lead to short-lived correlations between particles. The potential energy $U(i)$ of particle $i$ is determined additively:
\begin{align*}
U(i) = \sum_{j}U_{ij}.
\end{align*}
\item We say that a particle is {\em bound} to each of its neighbors with which it shares a bond. When this relation is extended to an equivalent relation, we call the equivalence classes {\em bound states}. Equivalently, these are the connected components in the graph generated by connecting particles with an edge if they share a bond.
\item The particles are immersed in a viscous heat bath with temperature $T = 1/\beta$ and collision coefficient $\gamma$. Viscosity introduces damping and fluctuations.
\item Adjacent particles can shield a particle from a portion of the viscosity of the bath. We take $\gamma$ to be the maximum viscosity experienced by a particle and $\gamma_{min}$ to be the minimum. For a lattice site $\mb{r}$, let $\epsilon(\mb{r})$ be the fraction of adjacent lattice sites that are unoccupied by particles. Then
\begin{align*}
\gamma(\mb{r}) = \gamma_{min} + (\gamma - \gamma_{min}){3\epsilon(\mb{r}) \over 2\epsilon(\mb{r}) + 1}
\end{align*}
gives the viscosity experienced by a particle at $\mb{r}$. The residual shielding of an additional adjacent particle increases as the number of adjacent particles increases. This shielding mechanism enables state-dependent differences in energy dissipation rates and reflects the general phenomena that dissipation involves the dispersal of energy. Particles will often move between adjacent lattice sites. During any such transition we set the viscosity experienced by the particle to the arithmetic mean of the viscosities at each site.
\item The particles are subject to a space-dependent external drive represented by a scalar field $F_{ext}$ coupled to particle mass. In general these drives could take any form, whether stochastic or deterministic, and interesting behavior might across a variety of choices. We choose to focus on physically-inspired drives that evoke the spectrum of solar radiation near Earth. In accordance with Planck's law, we take
\begin{align*}
F_{ext}((x, n)) \sim {x^{3} \over e^{\alpha x} - 1} \quad \textrm{and} \quad F_{ext}((x, y)) = \psi((x, y))F_{ext}((x, n)),
\end{align*}
where $\alpha$ is a constant tuning the strength of the drive and $\psi$ specifies how the drive varies in space. To emulate solar radiation near Earth, the constant $\alpha$ should be set such that the maximum of the distribution of $F_{ext}$ at the border is around an order of magnitude less than the strongest binding energies attained by particles in the system. Radiation cannot in general pass through particles unhindered. Therefore $\psi$ should vary across the lattice in a manner that reflects this screening effect. We could treat particles as perfect barriers or aborbers and set
\begin{align*}
\psi((x, y)) = \begin{cases}
0, &\textnormal{there exists a particle in } \{x\} \times (y, n], \\
1, &\textnormal{else}.
\end{cases}
\end{align*}
We could also instead consider particles as imperfect absorbers and assume an exponential falloff in intensity, setting
\begin{align*}
\psi((x, y)) = \mu_{1}e^{\mu_{2}\ell} \textnormal{ if there exist} \ell > 0 \textnormal{ particles in } \{x\} \times (y, n].
\end{align*}
The constants $\mu_{1}$ and $\mu_{2}$ should be set such that an appropriate normalization holds:
\begin{align*}
\sum_{\ell = 0}^{\infty}\mu_{1}e^{-\mu_{2}\ell} = 1,
\end{align*}
giving $\mu_{1} = 1 - e^{\mu_{2}}$.
\item Write $\mb{X} = (\mb{R}, \mb{E}, \mb{B}, \boldsymbol{\Sigma})$ for the state of the system, composed of the positions, energies, bond indicators, and bond stresses. The dynamics of the system are determined by elementary events $e: \mb{X} \leftarrow \mb{X}'$ that may occur in any given state. These events take place with state-dependent probabilities $p_{e}(\mb{X})$. There are three types of events:
\begin{itemize}
\item[--] Bound state hop $(A, \mb{e})$: bound state $A$ moves from $\mb{r}_{A}$ to $\mb{r}_{A} + \mb{e}$ for a shift vector $\mb{e}$, where $\mb{r}_{A}$ is the center of mass of $A$.
\item[--] Bond flip $(i, j)$: adjacent particles $i$ and $j$ change bond state from $b_{ij}$ to $\{1 - b_{ij}, b_{ij}\}$.
\item[--] Stress relaxation: bond stress diffuses between incident bonds.
\end{itemize}
Strong spatial particle correlations are mediated by bonding dynamics, obviating the need for separate particle-level sampling. We treat bond stress relaxation as occurring at much faster time scales than the other events. This enables us to perform the event deterministically between the other events, simplifying the dynamics and reducing sampling costs. The stochastic event types can mutually effect each other. As a result we run two separate phases of the sampling process, one for each event type. Furthermore, particles mutually interact, so movement events are not in general causally independent of one another. This requires us to factorize the set of bound states into independent subsets. We can find a partition into, in general, around 10 or fewer independent sets by iteratively choosing without replacement maximal independent sets using a sliding window algorithm.
\item ...
\end{itemize}

Draft model idea 7:
\begin{itemize}
\item Coordinate space is a cylindrical lattice $[n] \times [n]$ for $n \in \mathbb{N}$, with a periodic boundary on the first dimension and closed boundaries on the second; we write the points as vectors $\mb{r} = (x,y)$.\footnote{As previously, this is the default choice for concreteness and to mirror the geometry of Earth's surface. Most of the dynamics extend immediately to other locally finite lattices.}
\item The space is populated by $N$ indistinguishable particles, with positional Pauli exclusion amongst all particles. Each particle $i$ has a position $\mb{r}_{i}$ and internal energy $E_{i}$. We require
\begin{align*}
E_{i} \geq E_{0} > 0
\end{align*}
for a small constant $E_{0}$. The state of the particle system is
\begin{align*}
\mb{X} = (\mb{R}, \mb{E}, \mb{B}),
\end{align*}
where $\mb{R}$, $\mb{E}$, and $\mb{B}$ are the particle positions, internal energies, and bond indicators, respectively. % The viscous bath is assumed to be sufficiently overdamped that momentum is lost on a much faster scale than any modeled state transition.
\item We interpret $E_{i}$ as a finite internal thermal energy with constant effective heat capacity $C$. The effective temperature of particle $i$ is therefore
\begin{align*}
T_{i} = {E_{i} \over C}.
\end{align*}
Although we do not track internal entropy independently, this temperature relation implicitly determines the constitutive entropy
\begin{align*}
S(E) = C\log{E \over E_{0}},
\end{align*}
up to an additive constant, since $\partial S/\partial E = 1/T(E)$. This relation will only be used to determine the relative probabilities of transitions that exchange or transport internal energy. In particular, for $|\Delta E| \ll E$ we have
\begin{align*}
{S(E + \Delta E) - S(E) \over k_{B}} \approx {\Delta E \over k_{B}T(E)},
\end{align*}
recovering the usual Boltzmannian approximation without treating the particle as an infinite thermal reservoir.
\item Particles interact through a universal weak nearest-neighbor attraction and a stronger binary bonding interaction. Two distinct particles $i$ and $j$ have joint potential energy
\begin{align*}
U_{ij}(\mb{X}) = \begin{cases}
-1, &b_{ij} = 1, \\
-\eta, &b_{ij} = 0 \textrm{ and } |\mb{r}_{i} - \mb{r}_{j}| = 1, \\
0, &\textrm{else},
\end{cases}
\end{align*}
where $0 < \eta \ll 1$. A bond can only exist between adjacent particles, so $b_{ij}=0$ whenever $|\mb{r}_{i}-\mb{r}_{j}|>1$. The total configurational energy is
\begin{align*}
U(\mb{X}) = \sum_{i<j}U_{ij}(\mb{X}).
\end{align*}
Thus all close interactions are attractive and all energy barriers are configurational or geometric rather than being produced by particle polarity. The weak interaction is chosen so that ordinary contacts produce transient correlations and geometric trapping, while bonds produce much longer-lived structure.
\item We say that particles $i$ and $j$ are \emph{bound} if $b_{ij}=1$. When this relation is extended transitively, its equivalence classes are called \emph{bound states} or \emph{molecules}. Equivalently, the molecules are the connected components of the graph generated by joining bonded particles. Particle positions are fixed relative to one another within a molecule during a movement event. % Bond formation can merge molecules, while the breaking of a bridge bond can split one molecule into several molecules. In this way the number, sizes, geometries, and internal bond graphs of the effective interacting units are all dynamical.
\item The particles are immersed in a viscous heat bath with temperature $T_{b}$ and inverse temperature
\begin{align*}
\beta_{b} = {1 \over k_{B}T_{b}}.
\end{align*}
The bath has two distinct modeled effects. First, viscosity controls the kinetic frequency of molecule movements. Second, the bath exchanges heat with particle internal energies and may supply or absorb the energy required by structural transitions.
\item Geometry shields particles from the surrounding bath. For a particle at $\mb{r}_{i}$, let
\begin{align*}
\epsilon_{i}(\mb{X}) = {1 \over 4}\left|\{\mb{r}: |\mb{r}-\mb{r}_{i}|=1 \textrm{ and } \mb{r} \textrm{ is unoccupied}\}\right|
\end{align*}
be its exposed fraction. The local collision coefficient and thermal coupling are taken to depend linearly on exposure,
\begin{align*}
\gamma_{i}(\mb{X}) &= \gamma_{0}+\gamma_{1}\epsilon_{i}(\mb{X}), \\
\kappa_{i}(\mb{X}) &= \kappa_{0}+\kappa_{1}\epsilon_{i}(\mb{X}),
\end{align*}
where $\gamma_{0}>0$, $\gamma_{1}\geq 0$, and $0\leq\kappa_{0}\leq\kappa_{0}+\kappa_{1}\leq1$. The same exposure statistic is used for simplicity, but the mechanical and thermal parameters are independent. This makes both damping and heat loss state-dependent: exposed structures interact strongly with the bath, while shielded particles can retain energy for longer periods. The translational drag of a molecule $A$ is taken to be additive,
\begin{align*}
\Gamma_{A}^{\mathrm{tr}}(\mb{X}) = \sum_{i\in A}\gamma_{i}(\mb{X}).
\end{align*}
\item The external source is a stochastic scalar energy field incident from the top of the lattice. Fix an inverse source energy scale $\beta_{s}>0$ and let
\begin{align*}
f_{s}(q) = {15\beta_{s}^{4}\over\pi^{4}}{q^{3}\over e^{\beta_{s}q}-1}, \qquad q>0,
\end{align*}
be the normalized distribution of incident packet energies. At each irradiation phase and for every column $x$, independently sample
\begin{align*}
q_{x}\sim f_{s}.
\end{align*}
Let $i(x)$ be the topmost particle in column $x$. If the column is occupied, all incident energy is absorbed by this particle:
\begin{align*}
E_{i(x)} \leftarrow E_{i(x)} + q_{x}.
\end{align*}
If the column is empty, the packet leaves the system through the lower boundary. Thus the source is sampled independently of system state, while its absorption is state-dependent through shielding.
\item Internal energy is transported between adjacent particles by stochastic, energy-conserving pair exchanges. Fix a conduction quantum $\delta_{c}>0$. If particles $i$ and $j$ are adjacent, define the symmetric kinetic probability
\begin{align*}
\lambda_{ij} = \begin{cases}
\lambda_{b}, &b_{ij}=1, \\
\lambda_{c}, &b_{ij}=0,
\end{cases}
\end{align*}
where $0 \leq \lambda_{c}\leq\lambda_{b}\leq 1$, with a strict bonded advantage in the default model. Independently propose a direction $s\in\{-1,1\}$ with equal probability and the candidate update
\begin{align*}
E_{i}' &= E_{i}-s\delta_{c}, \\
E_{j}' &= E_{j}+s\delta_{c}.
\end{align*}
Thus $s=1$ transfers one quantum from $i$ to $j$, while $s=-1$ transfers it from $j$ to $i$. A proposal is illegal if either candidate energy is below $E_{0}$. Otherwise define
\begin{align*}
\mc{A}_{ij,s}={S(E_{i}-s\delta_{c})+S(E_{j}+s\delta_{c})-S(E_{i})-S(E_{j})\over k_{B}}
\end{align*}
and accept the exchange with probability
\begin{align*}
\lambda_{ij}a(\mc{A}_{ij,s}), \qquad a(\mc{A})=\min\{1,e^{\mc{A}}\}.
\end{align*}
The exchange conserves $E_{i}+E_{j}$ exactly and exchanges no heat with the bath. Because the directional proposal and $\lambda_{ij}$ are unchanged under reversal, its forward/reverse probability ratio is $e^{\mc{A}_{ij,s}}$. It therefore preserves the finite-reservoir equilibrium law rather than deterministically erasing its internal-energy fluctuations. In particular, with fixed geometry and no source, conduction preserves the product canonical distribution
\begin{align*}
\pi_{b}(\mb{E}) \propto \prod_{i}\mb{1}_{\{E_{i}\geq E_{0}\}}
\exp\left({S(E_{i})\over k_{B}}-\beta_{b}E_{i}\right),
\end{align*}
conditional on the energy classes reachable by the fixed quantum. The direct bath kernel introduced below preserves the same distribution, so their composition produces no artificial source-free heat current.

The nearest-neighbor edges of the square lattice are partitioned into four matching classes according to orientation and parity. Within a class no two edges share an endpoint, so all proposals are sampled in parallel. The four classes are applied sequentially in a newly randomized order, with energies recalculated after each class. Bonded particles conduct more frequently than particles in ordinary contact, so molecular topology controls transport kinetics without changing the equilibrium energy law. Energy remains localized on particles rather than being averaged over a molecule, allowing molecules to maintain fluctuating gradients and localized hot regions.
\item The dynamics are generated by a sequence of sampling phases rather than by tracking a global physical time. One complete application of all phases will be called a \emph{sweep}. The sweep index specifies only the order of updates; all transition probabilities and kinetic frequencies are dimensionless per-sweep parameters.
\item The stochastic structural dynamics are formulated using channel-resolved local detailed balance.\footnote{This is not a claim that the full driven system is in thermal equilibrium. It only constrains each bath-coupled microscopic channel; the continually sampled radiation source explicitly breaks global detailed balance.} Consider any reversible event $e:\mb{X}\to\mb{X}'$. Write
\begin{align*}
\Delta U_{e} &= U(\mb{X}')-U(\mb{X}), \\
\Delta E_{e} &= \sum_{i}(E_{i}'-E_{i}), \\
Q_{e} &= \Delta U_{e}+\Delta E_{e}, \\
\Delta S_{e} &= \sum_{i}\left(S(E_{i}')-S(E_{i})\right).
\end{align*}
The quantity $Q_{e}$ is the heat absorbed by the particle system from the bath during the event. The event affinity is
\begin{align*}
\mc{A}_{e} = {\Delta S_{e} \over k_{B}}-\beta_{b}Q_{e}.
\end{align*}
For every event and its reverse $\bar e$, choose a symmetric proposal probability and kinetic factor. Conditional on proposing $e$, accept it with the Metropolis function
\begin{align*}
a(\mc{A}) = \min\{1,e^{\mc{A}}\},
\end{align*}
which satisfies $a(\mc{A})/a(-\mc{A})=e^{\mc{A}}$. The resulting transition probabilities obey the required forward/reverse ratio. One could instead evaluate every candidate event and sample once from the resulting categorical distribution. We retain the proposal-and-acceptance formulation because it requires constructing the transformed footprint, contact changes, and energy update for only one candidate per molecule. The additional random draw is expected to be much cheaper than evaluating every translation, rotation, and energy channel, particularly for large molecules and large-scale parallel implementations.
\item Every structural event is allowed through two minimal energy channels. In the \emph{bath channel}, internal energies are unchanged, so
\begin{align*}
\Delta E_{e}=0, \qquad Q_{e}=\Delta U_{e}, \qquad \mc{A}_{e,B}=-\beta_{b}\Delta U_{e}.
\end{align*}
The bath may therefore power an arbitrarily costly uphill transition, but its probability is exponentially suppressed. In the \emph{internal channel}, configurational energy is exchanged locally with particle internal energies so that
\begin{align*}
\Delta E_{e}=-\Delta U_{e}, \qquad Q_{e}=0, \qquad \mc{A}_{e,I}={\Delta S_{e} \over k_{B}}.
\end{align*}
An uphill transition cools the relevant particles, while a downhill transition heats them. The internal channel is unavailable if it would send any particle energy below $E_{0}$. These are distinct microscopic pathways for the same structural change.
\item The first structural sampling phase consists of rigid molecule transformations. Fix a molecule $A$. A translation proposal is obtained by choosing
\begin{align*}
\mb{e}\in\{(1,0),(-1,0),(0,1),(0,-1)\}
\end{align*}
uniformly and mapping $\mb{r}_{i}\mapsto\mb{r}_{i}+\mb{e}$ for all $i\in A$. To avoid making molecular orientation an artificial permanent property, rotation proposals are also permitted. A rotation proposal is obtained by choosing a pivot particle $p\in A$ uniformly and rotating every particle of $A$ by either $\pi/2$ or $-\pi/2$ about $\mb{r}_{p}$ with equal probability. The use of a particle pivot ensures that the transformed positions remain on the lattice. A proposed transformation is illegal if it crosses a closed vertical boundary or maps a particle to a site occupied by a particle outside $A$.
\item Let $g$ be a legal molecule transformation and let $\mb{X}^{g}$ be the candidate state. Internal bonds and internal particle contacts are preserved by a rigid transformation. Consequently
\begin{align*}
\Delta U_{g}=U(\mb{X}^{g})-U(\mb{X})
\end{align*}
is determined entirely by weak contacts formed or broken between $A$ and particles outside $A$. For the internal channel, decompose this difference over changed contact pairs:
\begin{align*}
\Delta U_{g} = \sum_{\{i,j\}\in D_{g}}\Delta U_{ij},
\end{align*}
where $D_{g}$ is the set of cross-molecule pairs whose contact state changes. The local energy update is
\begin{align*}
E_{i}' &\leftarrow E_{i}'-{\Delta U_{ij}\over 2}, \\
E_{j}' &\leftarrow E_{j}'-{\Delta U_{ij}\over 2}
\end{align*}
for every $\{i,j\}\in D_{g}$, beginning from $\mb{E}'=\mb{E}$. Thus a newly formed weak contact deposits $\eta/2$ into each endpoint, while a broken weak contact removes $\eta/2$ from each endpoint. This assignment is local, symmetric, and satisfies $\Delta E_{g}=-\Delta U_{g}$ exactly.
\item The kinetic mobility of a proposed molecule transformation is determined by molecular temperature and viscous drag. Let
\begin{align*}
T_{A}(\mb{X}) = {1 \over C|A|}\sum_{i\in A}E_{i}
\end{align*}
be the mean molecular temperature. For an event $(g,\alpha):\mb{X}\to\mb{X}^{g,\alpha}$ through channel $\alpha\in\{B,I\}$ define the symmetric temperature
\begin{align*}
T_{g,\alpha} = {T_{A}(\mb{X})+T_{A}(\mb{X}^{g,\alpha}) \over 2}.
\end{align*}
For a translation, define the symmetric drag
\begin{align*}
\Gamma_{g}^{\mathrm{tr}} = {\Gamma_{A}^{\mathrm{tr}}(\mb{X})+\Gamma_{A}^{\mathrm{tr}}(\mb{X}^{g,\alpha}) \over 2}
\end{align*}
and kinetic factor
\begin{align*}
m_{g,\alpha}^{\mathrm{tr}} = \min\left\{1,\nu_{\mathrm{tr}}{T_{g,\alpha} \over T_{b}}{\Gamma_{0}^{\mathrm{tr}} \over \Gamma_{g}^{\mathrm{tr}}}\right\}.
\end{align*}
For a rotation about pivot particle $p\in A$, define the rotational drag
\begin{align*}
\Gamma_{A,p}^{\mathrm{rot}}(\mb{X}) = \sum_{i\in A}\gamma_{i}(\mb{X})\|\mb{r}_{i}-\mb{r}_{p}\|_{2}^{2}.
\end{align*}
Symmetrizing between the two states gives
\begin{align*}
\Gamma_{g,p}^{\mathrm{rot}} = {\Gamma_{A,p}^{\mathrm{rot}}(\mb{X})+\Gamma_{A,p}^{\mathrm{rot}}(\mb{X}^{g,\alpha})\over 2}.
\end{align*}
The rotational kinetic factor is
\begin{align*}
m_{g,\alpha}^{\mathrm{rot}} = \min\left\{1,\nu_{\mathrm{rot}}{T_{g,\alpha} \over T_{b}}{\Gamma_{0}^{\mathrm{rot}} \over \Gamma_{g,p}^{\mathrm{rot}}}\right\}.
\end{align*}
This form follows from viscous dissipation. Under angular velocity $\omega$, particle $i$ moves with speed $\omega\|\mb{r}_{i}-\mb{r}_{p}\|_{2}$, so the total dissipated power is proportional to
\begin{align*}
\omega^{2}\sum_{i\in A}\gamma_{i}\|\mb{r}_{i}-\mb{r}_{p}\|_{2}^{2}=\omega^{2}\Gamma_{A,p}^{\mathrm{rot}}.
\end{align*}
Rotational mobility is therefore inversely proportional to this drag. The fixed quarter-turn angle is absorbed into $\nu_{\mathrm{rot}}$. Rotation proposals are omitted for monomers. Let $m_{g,\alpha}$ denote the appropriate translational or rotational factor. Since the temperatures and drags are symmetrized between an event and its reverse, these factors change transition speeds without changing the thermodynamic forward/reverse ratio.
\item The molecule transformation phase is sampled as follows. For each selected molecule $A$:
\begin{itemize}
\item[--] Propose a translation with probability $\zeta$ or a rotation with probability $1-\zeta$, using the proposal distributions above.
\item[--] Independently choose the bath channel with probability $\chi_{B}$ or the internal channel with probability $\chi_{I}=1-\chi_{B}$.
\item[--] Reject an illegal geometric proposal or an internal-channel proposal producing an energy below $E_{0}$.
\item[--] Otherwise accept the event with probability
\begin{align*}
m_{g,\alpha}a(\mc{A}_{g,\alpha}),
\end{align*}
where $\alpha\in\{B,I\}$ is the sampled channel and $m_{g,\alpha}$ is the appropriate translational or rotational kinetic factor.
\end{itemize}
If the proposal is rejected, the molecule remains in place.
\item Molecule transformations are not mutually independent because transformed molecules may collide or change common contact neighborhoods. We use an optimistic parallel scheduler rather than partitioning molecules according to every transformation they could possibly make. At the beginning of the phase all molecules are unresolved. From a common state snapshot, every unresolved molecule independently samples one transformation, one energy channel, and its Metropolis acceptance decision in parallel. Illegal and rejected proposals are resolved as stays. For every provisionally accepted proposal $g$, let $\mc{C}(g)$ be the union of its source and destination footprints and their nearest-neighbor shells. Two proposals conflict when their claim sets intersect. Assign the active proposals independent random priorities and repeatedly select proposals that have maximal priority on every claimed site, removing selected proposals and their conflicts until a maximal independent set is obtained. Apply this set simultaneously and mark its molecules as resolved. Provisionally accepted proposals excluded by a conflict remain unresolved and are resampled from the updated state. The procedure continues until every molecule is resolved or a fixed iteration limit is reached, in which case the remaining molecules stay in place. Thus every molecule undergoes at most one transformation in a phase, while conflicts are computed only for transformations that were actually sampled and provisionally accepted. This is expected to expose substantially more parallelism than a conservative precomputed partition. The state-dependent conflict filter can perturb the exact global equilibrium kernel when accepted proposals conflict, so the source-free dynamics should be validated against a conservative independent-set implementation; the two schemes coincide whenever the active proposals are conflict-free.
\item The second structural sampling phase consists of bond flips. Fix two adjacent particles $i$ and $j$ and let $e=(i,j)$ be the event
\begin{align*}
b_{ij}'=1-b_{ij}.
\end{align*}
Its configurational energy difference is
\begin{align*}
\Delta U_{e}=(2b_{ij}-1)(1-\eta).
\end{align*}
Thus bond formation releases $1-\eta$ units of configurational energy and bond breaking requires the same amount. In the internal channel this energy is divided symmetrically between the endpoints:
\begin{align*}
E_{i}' = E_{i}-{\Delta U_{e}\over 2}, \qquad E_{j}' = E_{j}-{\Delta U_{e}\over 2}.
\end{align*}
The bath channel leaves both internal energies unchanged. Bond formation between particles in different molecules merges their components immediately, while a bond break is followed by a connected-component update that may split the original molecule.
\item Bond flips are activated events. Their energy difference determines the forward/reverse bias, while a separate symmetric kinetic prefactor determines how often the barrier is crossed. To supply a minimal endogenous mechanism by which existing structure can alter the production of future structure, we include geometric catalysis. For a candidate lattice edge $e$, let $\Phi_{e}(\mb{X})$ be the number of bonded edges parallel to $e$ at perpendicular lattice distance one. The candidate edge itself is excluded. Define
\begin{align*}
c_{e}(\mb{X}) = \min\left\{c_{\max},\exp(\lambda_{\mathrm{cat}}\Phi_{e}(\mb{X}))\right\},
\end{align*}
where $\lambda_{\mathrm{cat}}\geq 0$. Because $\Phi_{e}$ does not depend on the state of the candidate bond, it is unchanged by the flip and multiplies the forward and reverse rates equally. It therefore lowers a local activation barrier without supplying energy or changing the required detailed-balance ratio. Setting $\lambda_{\mathrm{cat}}=0$ recovers the uncatalyzed model.
\item Let
\begin{align*}
T_{ij}(\mb{X})={E_{i}+E_{j}\over 2C}
\end{align*}
be the mean endpoint temperature. For a candidate bond flip $(e,\alpha):\mb{X}\to\mb{X}^{e,\alpha}$ define
\begin{align*}
\overline{T}_{e,\alpha}={T_{ij}(\mb{X})+T_{ij}(\mb{X}^{e,\alpha})\over 2}
\end{align*}
and the symmetric kinetic factor
\begin{align*}
m_{e,\alpha}^{\mathrm{bond}}=\min\left\{1,\nu_{\mathrm{bond}}{\overline{T}_{e,\alpha}\over T_{b}}c_{e}(\mb{X})\right\}.
\end{align*}
The bond phase is sampled by independently choosing the bath or internal channel with probabilities $\chi_{B}$ and $\chi_{I}$ and accepting the flip with probability
\begin{align*}
m_{e,\alpha}^{\mathrm{bond}}a(\mc{A}_{e,\alpha}).
\end{align*}
The base frequency $\nu_{\mathrm{bond}}$ represents the uncatalyzed activation barrier. Hot endpoints react more frequently, while the affinity determines whether the selected direction is thermodynamically favored.
\item Candidate bond edges are not all independent because simultaneous flips may share endpoints or change one another's catalytic neighborhoods. There is a fixed four-color edge partition that permits parallel sampling within each class of candidate bonds. Sequential sampling of each class can be randomly permuted to avoid introducing update-order artifacts into the proposed templating mechanism.
\item Direct heat exchange between particles and the bath forms a separate sampling phase. This is necessary because a stationary or strongly bound particle should still be capable of heating and cooling. Fix a bath-exchange quantum $\delta_{b}>0$. For each particle $i$, independently propose $s\in\{-1,1\}$ with equal probability and the candidate update
\begin{align*}
E_{i}'=E_{i}+s\delta_{b}.
\end{align*}
The negative update is illegal if $E_{i}'<E_{0}$. Otherwise its affinity is
\begin{align*}
\mc{A}_{i,s}={S(E_{i}+s\delta_{b})-S(E_{i})\over k_{B}}-\beta_{b}s\delta_{b}.
\end{align*}
The update is accepted with probability
\begin{align*}
\kappa_{i}(\mb{X})a(\mc{A}_{i,s}),
\end{align*}
where the constants defining $\kappa_{i}$ are chosen so that $0\leq\kappa_{i}\leq1$. This kernel preserves $\pi_{b}$ when the geometry is fixed; its modal energy is near $E=CT_{b}$. Exposed particles exchange heat rapidly, while shielded particles retain internal energy. Exposure changes the relaxation speed but not the stationary energy law. All particle proposals are conditionally independent in this phase because the geometry is held fixed.
\item One complete sweep consists of the following phases:
\begin{itemize}
\item[1.] \emph{Irradiation}: sample one incident packet for every column and add it to the topmost particle.
\item[2.] \emph{Conduction}: sample one signed pair exchange over each of the four nearest-neighbor matching classes in randomized class order.
\item[3.] \emph{Molecule transformations}: sample rigid translations and rotations in optimistic parallel batches, filtering provisionally accepted proposals to maximal independent sets.
\item[4.] \emph{Bond dynamics}: sample bond flips over independent edge classes, updating the molecule decomposition after each class.
\item[5.] \emph{Conduction}: optionally repeat the stochastic conduction phase with independent proposals so energy released by structural changes can begin propagating during the same sweep. Repeating this phase increases the effective per-sweep conduction frequency and should be calibrated accordingly.
\item[6.] \emph{Bath exchange}: sample direct particle--bath heat exchanges.
\end{itemize}
The relative ordering of the molecule and bond phases can be randomly reversed at the beginning of each sweep to reduce systematic advantages from a fixed update order. % The radiation source acts exactly once per sweep regardless of whether the system has equilibrated. In particular, global ground states have no privileged role in the dynamics.
\item The proposed model contains only one particle type, one scalar particle energy, one binary bond type, nearest-neighbor interactions, a single external source, and a single bath. Nevertheless, the connected-component construction gives an enormous dynamical space of effective entities. Molecules can capture radiation according to their geometry, transport it through their bond networks, retain it through shielding, use it to cross configurational barriers, merge, fragment, and alter the kinetics of nearby bond formation through geometric catalysis. The intended recursive feedback is
\begin{align*}
\textrm{structure} &\longrightarrow \textrm{energy capture, transport, and dissipation} \\
&\longrightarrow \textrm{state-dependent structural transitions} \\
&\longrightarrow \textrm{new structure}.
\end{align*}
The catalytic factor is the one addition included specifically for open-endedness. Without it, structures can persist and control energy flow but have no direct mechanism for altering the production of related future structures. Parallel-edge catalysis supplies a minimal form of templating using only existing geometry and the bond bit. Weak nonbonding attraction can hold a prospective product near a template, while molecule movement can later separate them. Reversible bonding ensures continual matter recycling rather than irreversible conversion into a single bonded aggregate.
\item This system cannot exhibit literally unbounded novelty on a fixed finite lattice. The scale hypothesis is instead that, for sufficiently large $n$ and $N$, the reachable space of molecule geometries, bond graphs, internal energy distributions, weak complexes, and catalytic interactions may be large enough to produce practical open-endedness on all accessible simulation scales. If the dynamics instead collapse to a small family of top-layer sheets or a permanently bonded gel, the first changes to consider are, in order: partial transmission or saturation of incident radiation, one additional particle or bond recognition bit, and more general local geometric catalysis. Bond stress should only be reconsidered if the lack of elastic or mechanically transmitted memory appears to be the principal limitation.
\end{itemize}


\begin{comment}
Work done on the system is stored as kinetic or potential energy in particles and bonds until it is released into the viscous heat bath via damping effects. By default a particle travels according to its momentum. Sometimes heat is borrowed from the bath to alter the trajectory of the particle. We determine the net effect of these processes by tracking the alignment between the velocity of the particle and its change in position, an analogue for the classical expression for work as force applied along a displacement. Thus the work dissipated from the particle system during an event $e$ is
\begin{align*}
W_{e}(\mb{X}) = {\mb{P}' \over \mb{M}} \cdot (\mb{R}' - \mb{R}),
\end{align*}
where $\mb{P}'$ is the momenta after $e$ and $\mb{M}$ is the vector of particle masses.
\end{comment}

\begin{comment}
\item Write $\mb{R}$, $\mb{P}$, and $\mb{B}$ for the particle positions, momenta, and bond states, respectively, where the latter consists of the bond indicators $\{b_{ij}\}$ and stresses $\{\boldsymbol{\sigma}_{ij}\}$. These constitute a point in the state space of the particle system, excluding the values of the external and Brownian forces. The dynamics of the system are determined by the probability distribution for the next state, $P(\mb{R}', \mb{P}', \mb{B}')$, where $\mb{R}', \mb{P}', \mb{B}'$ are the next positions, momenta, and bond states. There are at most five possible positions for the next state of a given particle, namely, the current position plus each of the four neighboring positions. The momentum is determined by the current state and chosen next position and bond state, so $P(\mb{R}', \mb{P}', \mb{B}') = P(\mb{R}', \mb{B}')$. The probability distribution for the next state is then given by the master equation
\begin{align*}
P(\mb{R}', \mb{B}') = \sum_{\mb{S}, \mb{A}}P(\mb{R}', \mb{B}' \mid \mb{S}, \mb{Q}, \mb{A})P(\mb{R} = \mb{S}, \mb{P} = \mb{Q}, \mb{B} = \mb{A})
\end{align*}
which reduces to
\begin{align*}
P(\mb{R}', \mb{B}') = P(\mb{R}', \mb{B}' \mid \mb{R}, \mb{P}, \mb{B})
\end{align*}
upon fixing $\mb{R}$, $\mb{P}$, and $\mb{B}$. So we must determine the conditional next state distribution $P(\mb{R}', \mb{B}' \mid \mb{R}, \mb{P}, \mb{B})$.
\item We begin by employing Gibbs sampling to separate the sampling of particle positions from their bond states. We first sample from the distribution $P(\mb{B}' \mid \mb{R}, \mb{P}, \mb{B})$ and then use the obtained samples $\mb{B}'$ to sample from $P(\mb{R}' \mid \mb{R}, \mb{P}, \mb{B}, \mb{B}')$. This has the advantage of reducing the number of distribution components in the sampling process and allows us to take advantage of the fact that bond states can be sampled independently.
\item Particles mutually interact, so the position variables are not in general causally independent of one another. To make sampling from $P(\mb{R}' \mid \mb{R}, \mb{P}, \mb{B}, \mb{B}')$ tractable this requires us to factorize the conditional distribution into independent pieces. Given that particles only have access to the five nearest lattice sites at each update, and both interaction mechanisms only extend out to these same positions, particles separated by a distance of four or more are pairwise independent. It follows that we may partition the particles into eight independent sets $I_{1}, \ldots, I_{8}$. Then Gibbs sampling over the partition allows us to sample from $P(\mb{R}' \mid \mb{R}, \mb{P}, \mb{B}, \mb{B}')$. When sampling the appropriate conditional distribution of an independent set $I_{\ell}$ of particles we can sample the next position of each particle in $I_{\ell}$ independently. This reduces the problem to sampling the next position of a single particle.
\item While this method avoids treating correlated particles as independent, it does not account for all of the effects of correlation. Highly correlated particles cannot be correctly sampled separately---the marginal distribution collapses in a way that does not reflect the global distribution. By construction this occurs most prominently when particles are in a bound state. To correct for this we run between the particle-level updates a third phase of blocked Gibbs sampling on bound states. For the purposes of efficient sampling we treat particles in a bound state as perfectly correlated during this phase, fixing their positions relative to one another as well as fixing all bond states. This reduces the state variables to the position of the center of mass of a bound state, as well as the orientation about its center of mass. To maintain simplicity and tractability, we elect to fix the orientation. This allows us to treat the bound state in the same manner as a particle, with five candidate positions for the next state of the center of mass. It is also enabled by the fact that rotation would have to occur in $\pi/2$ radian increments, a usually larger scale change than a translation by a single unit distance. This suggests that such changes are more rare and thus can be neglected.

\item We conclude the bond dynamics phase by sampling from
the distribution $P(\mb{B}' \mid \mb{R}, \mb{P}, \mb{B})$ over the bond states of all adjacent particles. Fix two adjacent particles $i$ and $j$ and a bit $b \in \{0, 1\}$. Let $\mb{B}'$ be the bond states obtained from $\mb{B}$ by setting $b_{ij} = b$ and $\mb{P}'$ be obtained from $\mb{P}$ by assigning particle $i$ and $j$ momenta
\begin{align*}
\mb{p}'_{i} = \mb{p}_{i} - {\nabla_{ij}U \over 2}\Delta t_{i} \quad \textrm{and} \quad \mb{p}'_{j} = \mb{p}_{j} + {\nabla_{ij}U \over 2}\Delta t_{j},
\end{align*}
respectively, where
\begin{align*}
\Delta t_{i} = \min\left\{{m_{i}\tau \over |\mb{p}_{i}|}, \sqrt{2m_{i}\tau \over |\nabla_{ij}U|}\right\} \quad \textrm{and} \quad \Delta t_{j} = \min\left\{{m_{j}\tau \over |\mb{p}_{j}|}, \sqrt{2m_{j}\tau \over |\nabla_{ij}U|}\right\}
\end{align*}
for a constant $\tau$ determining the impulse scale of bonding events relative to particle kinetics. Each particle is influenced symmetrically and in opposite directions by the bond state change. The heat released into the bath in the transition from $\mb{B}$ to $\mb{B}'$ is
\begin{align*}
\Delta Q_{\mb{B}\mb{B}'} &= -\Delta E_{\mb{B}\mb{B}'} \\
&= -\Delta K_{\mb{B}\mb{B}'} - \Delta U_{\mb{B}\mb{B}'} \\
&= - {1 \over 2m_{i}}(|\mb{p}'_{i}|^{2} - |\mb{p}_{i}|^{2}) - {1 \over 2m_{j}}(|\mb{p}'_{j}|^{2} - |\mb{p}_{j}|^{2}) - (U(\mb{B}') - U(\mb{B})) \\
&= - {1 \over 2m_{i}}(|\mb{p}'_{i}|^{2} - |\mb{p}_{i}|^{2}) - {1 \over 2m_{j}}(|\mb{p}'_{j}|^{2} - |\mb{p}_{j}|^{2}) \\
&\qquad\qquad - (b - b_{ij})\left((1 - \eta)q_{i}q_{j} + {1 \over 2}(k_{1}\sigma_{1,ij}^{2} + k_{2}\sigma_{2,ij}^{2})\right),
\end{align*}
where $\Delta E_{\mb{B}\mb{B}'}$, $\Delta K_{\mb{B}\mb{B}'}$, and $\Delta U_{\mb{B}\mb{B}'}$ are the energy, kinetic energy, and potential energy differences between $\mb{B}'$ and $\mb{B}$. As previously missing bond variables default to 0. Following Boltzmann thermodynamics we take
\begin{align*}
P(\mb{B}' \mid \mb{R}, \mb{P}, \mb{B}) = {1 \over Z}e^{\beta\Delta Q_{\mb{B}\mb{B}'}},
\end{align*}
where $Z$ is the partition function
\begin{align*}
Z = \sum_{\mb{B}'}e^{\beta\Delta Q_{\mb{B}\mb{B}'}} = \sum_{b}e^{\beta\Delta Q_{\mb{B}\mb{B}'}}.
\end{align*}
\item The second phase is the particle-level update phase where we sample from the distribution $P(\mb{R}' \mid \mb{R}, \mb{P}, \mb{B}')$ over the next position of each particle. Fix a particle $i$ of charge $q$ and mass $m$ and fix a shift vector $\mb{e} \in \{(0, 0), (1, 0), (-1, 0), (0, 1), (0, -1)\}$. Write
\begin{align*}
\mb{f}_{i} = q\mb{F}_{charge}(\mb{r}_{i}) + m\mb{F}_{mass}(\mb{r}_{i}) - \nabla_{\mb{e}}U(\mb{r}_{i}; i)
\end{align*}
for the net non-Brownian force experienced by particle $i$. Let $\mb{R}'$ be the positions matching $\mb{R}$ except particle $i$ is in position $\mb{r}_{i}' = \mb{r}_{i} + \mb{e}$ with momentum
\begin{align}
\label{mom_update_eq}
\mb{p}_{i} = (1 - \gamma\Delta t)\mb{p}_{i} + \mb{f}_{rand}\sqrt{\Delta t} + \mb{f}_{i}\Delta t,
\end{align}
where
\begin{align*}
\Delta t = \min\left\{{mD \over |\mb{p}_{i}|}, \sqrt{mD \over |\mb{f}_{i} - \gamma\mb{p}_{i}|}, \left({mD \over |\mb{f}_{rand}|}\right)^{2/3}\right\}.
\end{align*}
This gives a time scale approximately matching the expected position transition rates. The terms of (\ref{mom_update_eq}) track the change in momentum due to the viscosity of the heat bath and the various forces. Work done on the system is stored as kinetic or potential energy in particles and bonds until it is released into the viscous heat bath via damping effects. By default a particle travels according to its momentum. Sometimes heat is borrowed from the bath to alter the trajectory of the particle. We determine the net effect of these processes by tracking the alignment between the velocity of the particle and its change in position, an analogue for the classical expression for work as force applied along a displacement. Thus the work dissipated from the particle system in the transition from $\mb{R}$ to $\mb{R}'$ is
\begin{align*}
W_{\mb{R}\mb{R}'} = {1 \over m}(\mb{p}_{i}' \cdot \mb{e}).
\end{align*}
Then the heat $\Delta Q_{\mb{R}\mb{R}'}$ released into the bath in the transition between $\mb{R}'$ and $\mb{R}$ is
\begin{align*}
\Delta Q_{\mb{R}\mb{R}'} &= W_{\mb{R}\mb{R}'} - \Delta K_{\mb{R}\mb{R}'} - \Delta U_{\mb{R}\mb{R}'} \\
&= {1 \over m}(\mb{p}_{i}' \cdot \mb{e}) - {1 \over 2m}(|\mb{p}_{i}'|^{2} - |\mb{p}_{i}|^{2}) - (U(\mb{R}') - U(\mb{R})) \\
&= {1 \over m}(\mb{p}_{i}' \cdot \mb{e}) - {1 \over 2m}(|\mb{p}_{i}'|^{2} - |\mb{p}_{i}|^{2}) - (U(\mb{r}_{i}'; i) - U(\mb{r}_{i}; i)).
\end{align*}
Following non-equilibrium Boltzmann thermodynamics\footnote{Cite England, etc.} we take
\begin{align*}
P(\mb{R}' \mid \mb{R}, \mb{P}, \mb{B}') = {1 \over Z}e^{\beta\Delta Q_{\mb{R}\mb{R}'}},
\end{align*}
where
\begin{align*}
Z = \sum_{\mb{R}'}e^{\beta\Delta Q_{\mb{R}\mb{R}'}} = \sum_{\mb{e}}e^{\beta\Delta Q_{\mb{R}\mb{R}'}}.
\end{align*}
\end{comment}


\section{Notes}

\textit{Dissipative Structures, Organisms, and Evolution}
\begin{itemize}
\item External factors introduce or bias fluctuations that change the phenotype or macroscopic features of a non-equilibrium dissipative system; fluctuation-amplification
\item Autocatalysis or self-production is when a thermodynamic process is driven by the products of that process. Manifests in nearly all self-organizing phenomena.
\end{itemize}




\end{document}
